\documentclass[10pt,a4paper]{article}

% ----------------- Paquetes -------------------%
\usepackage[T1]{fontenc}
\usepackage[utf8]{inputenc}
\usepackage[a4paper,margin=2.5cm]{geometry}
\usepackage{mathptmx}             
\usepackage{changepage} 
\usepackage{setspace}
\usepackage[spanish]{babel}
\usepackage[labelfont=bf,labelsep=space]{caption}
\usepackage{amssymb}
\usepackage{amsmath}
\usepackage{ebgaramond}
\usepackage{placeins}
\usepackage{etoolbox}
\usepackage{setspace}
\usepackage{parskip} 
\usepackage{tcolorbox}
\usepackage{needspace}
\usepackage{xparse}
\usepackage{ocgx2}
\usepackage{hyperref}
\usepackage{hypcap}
\usepackage{soul}\sethlcolor{yellow}% resaltado amarillo: \hl{texto} (Panel TCEE, ⌃H)


%%%%%%%%%%%%%%%%%%%%%%%%%%%%%%%%%%%%%%%%%%%%%%%%%%%%%%%%%%%%%%%%
% --- Fija primero tus valores globales "buenos" ---
\setstretch{1.2}                 % Interlineado global
\setlength{\parskip}{0.7\baselineskip}
\setlength{\parindent}{0pt}
\newlength{\GlobalParskip}
\setlength{\GlobalParskip}{\parskip}

\newcounter{eqnum}[subsection]
\renewcommand{\theeqnum}{\arabic{eqnum}}
\newcounter{alphaitem}
\newcounter{numitem}

%% ------------------- Recuadros----------------------%%
\newcommand{\puntosclave}[1]{%
  \begin{tcolorbox}[
    colframe=black,
    title={Puntos clave},
    before upper={\setlength{\parindent}{0.5cm}\setlength{\parskip}{0.5\baselineskip}}
  ]
  #1
  \end{tcolorbox}
}

\newcommand{\modificaciones}[1]{
  \begin{tcolorbox}[
    colback=white,
    colframe=red,
    title={Modificaciones a realizar},
    before upper={\setlength{\parindent}{0.5cm}\setlength{\parskip}{0.5\baselineskip}}
  ]
  #1
  \end{tcolorbox}
}

%% ------------------- CITA ----------------------%%
% \begin{cita}[Autor][Año][Obra] texto \end{cita}: cita sangrada y, debajo a la derecha, «AUTOR (Año), Obra». Los tres datos son opcionales.
% En el Panel TCEE: escribir \cita e Intro (el cursor va al texto; Tab pasa a Autor, Año y Obra).
\NewDocumentEnvironment{cita}{ O{} O{} O{} +b }{%
  \begin{adjustwidth}{2cm}{0pt}
  \raggedright
  #4\par
  \raggedleft
  \if\relax\detokenize{#1#2#3}\relax
    % sin firma
  \else
    #1%
    \if\relax\detokenize{#2}\relax\else\if\relax\detokenize{#1}\relax\else\space\fi(#2)\fi
    \if\relax\detokenize{#3}\relax\else\if\relax\detokenize{#1#2}\relax\else, \fi\textit{#3}\fi
  \fi
  \end{adjustwidth}
}{}



%% ------------------- Listas----------------------%%
    % --- Lista alfabética ---
    \newenvironment{la}
      {\begin{list}{\alph{alphaitem})}{%
          \usecounter{alphaitem}%
          \setlength{\leftmargin}{1cm}%
          \setlength{\parskip}{3pt}% 
          \setlength{\itemsep}{0pt}% ← espacio entre ítems
          \setlength{\parsep}{0pt}%  ← párrafos dentro de un ítem
      }}
      {\end{list}}
      
    % --- Lista numérica ---
    \newenvironment{lnum}
      {\begin{list}{\arabic{numitem}.}{%
          \usecounter{numitem}%
          \setlength{\leftmargin}{1cm}%
          \setlength{\parskip}{3pt}% 
          \setlength{\itemsep}{0pt}% ← espacio entre ítems
          \setlength{\parsep}{0pt}%  ← párrafos dentro de un ítem
      }}
      {\end{list}}
      
% ------------------- Imágenes----------------------%
    \graphicspath{{Img/}}% Rutas por defecto 
    % --- Imagen adaptativa ---
    % Registros de dimensión definidos una sola vez
    \newdimen\imgcap
    \newdimen\imgmin
    \newdimen\imgmax
    \newdimen\imgremain
    \newdimen\imgh
    \newdimen\imgneed
    
    \newcommand{\imagenfit}[3]{%
      \addvspace{10pt}% separación antes
      \par\begingroup
        % Parámetros de composición (asignaciones, no \newdimen)
        \imgcap=1\baselineskip            % alto estimado de título+fuente
        \imgmin=0.15\textheight
        \imgmax=0.15\textheight
        \imgremain=\pagegoal
        \ifdim\pagegoal=\maxdimen \imgremain=0pt\fi
        \advance\imgremain by -\pagetotal
        \advance\imgremain by -\imgcap
        % Decisión de altura destino
        \ifdim\imgremain<\imgmin
          \newpage
          \imgh=0.20\textheight
        \else
          \imgh=\imgremain
          \ifdim\imgh>\imgmax \imgh=\imgmax \fi
        \fi
        % Reservar espacio para evitar cortes del bloque completo
        \imgneed=\imgh \advance\imgneed by \imgcap
        \Needspace{\imgneed}
        % Cuerpo
        \noindent\begin{minipage}{\textwidth}
          \centering
          \captionof{figure}{\textemdash\ #2}\par\vspace{0.3em}% título arriba
          \includegraphics[width=\linewidth,height=\imgh,keepaspectratio]{\detokenize{#1}}\par
          \vspace{0.3em}\small\textit{Fuente: }#3% fuente abajo
        \end{minipage}%
      \endgroup
      \par\addvspace{10pt}%
    }

%------------ Bloque de RECUADRO ESPECIAL -------------%
    \newenvironment{specialblock}[1]{%
      \begin{quote} % crea sangría a izquierda y derecha
      \small % texto más pequeño
      \noindent\textbf{#1}\\[-0.3em] % título en negrita
      \rule{\linewidth}{0.4pt}\\[0.6em]}{
      \end{quote}}
      

%-------------------------- Author Font -----------------------%
    \newcommand{\authorfont}[1]{{\fontfamily{EBGaramond-TLF}\selectfont #1}}

%-------------------------- Anexos -----------------------%
    \makeatletter
    \newcounter{anexo}
    \renewcommand{\theanexo}{\Roman{anexo}}
    \newcommand{\anexo}[1]{%
      \clearpage
      \phantomsection
      \refstepcounter{anexo}%
      \noindent\textbf{Anexo \theanexo.- }#1\par\medskip
      \addcontentsline{stc}{section}{Anexo \theanexo.- #1}%
    }
    \makeatother

    %-------------------------- Lista de figuras (personal) -----------------------%
\makeatletter
\newcommand{\MyListOfFigures}{%
  \clearpage
  \phantomsection
  {\centering\bfseries\Large Índice de figuras\par}%
  \medskip
  % Entrada en nuestro índice de contenido (stc)
  \addcontentsline{stc}{section}{Índice de figuras}%
  % Recuperación del listado (sin cabecera propia)
  \begingroup
    \parindent=0pt\parskip=2pt
    \@starttoc{lof}%
  \endgroup
}
\makeatother

%-------------------------- Bibliografía -----------------------%
\makeatletter
% Contador para las referencias
\newcounter{myref}
% Cabecera de la bibliografía, entra en el índice 'stc'
\newcommand{\MyBibliography}{%
  \clearpage
  \phantomsection
  {\centering\bfseries\Large Bibliografía y referencias\par}%
  \medskip
  \addcontentsline{stc}{section}{Bibliografía y referencias}%
  \setcounter{myref}{0}%
}
\newcommand{\MyBibItem}[4]{%
  \refstepcounter{myref}%
  \noindent[\themyref]\ #1.\ \emph{#2}. #3. #4\par
}
\makeatother

%--------- Establecer cuatro niveles de índice --------%
    \setcounter{tocdepth}{4}   
    \setcounter{secnumdepth}{4}% numera hasta \paragraph

%%%%%%%%%%%%%%%%%%%%%%%%%%%%%%%%

\makeatletter

    %--------- Interlineado Global --------%
    \edef\GlobalBaseLineStretch{\baselinestretch}
    
    %--------- Espaciado de seccinones --------%
    \renewcommand\section{\@startsection{section}
      {1}{\z@}
      {2.5ex \@plus 1ex \@minus .2ex} % espacio antes
      {1.5ex \@plus .2ex}             % espacio después
      {\normalfont\Large\bfseries}    % formato del título
    }
    \renewcommand\subsection{\@startsection{subsection}{2}{\z@}
      {3ex \@plus 0.8ex \@minus .2ex}
      {1.5ex \@plus .2ex}
      {\normalfont\large\bfseries}
    }
    \renewcommand\subsubsection{\@startsection{subsubsection}{3}{\z@}
      {3ex \@plus 0.5ex \@minus .2ex}
      {1ex \@plus .2ex}
      {\normalfont\normalsize\bfseries}
    }
    \renewcommand\paragraph{\@startsection{paragraph}{4}{\z@}%
    {-2ex \@plus -0.5ex \@minus -.2ex}% espacio antes
    {0.8ex \@plus .2ex}% espacio después
    {\normalfont\normalsize\itshape}}% ← cursiva en lugar de negrita

    %--------- Ecuaciones --------%   
    % Variables globales para almacenar títulos y notas
    \gdef\leq@current@section{}%
    \gdef\leq@current@subsection{}%
    \gdef\leq@last@printed@section{}%
    \gdef\leq@last@printed@subsection{}%
    
    % Comando para añadir nota (invisible en el cuerpo)
    \newcommand{\eqnote}[1]{%
      \addcontentsline{leq}{leqnote}{#1}%
    }
    
    % Comando para ecuaciones
    \newcommand{\eqblock}[2]{%
      % Verificar si necesitamos imprimir el título de sección
      \ifx\leq@current@section\leq@last@printed@section\else
        \addcontentsline{leq}{leqsection}{\leq@current@section}%
        \global\let\leq@last@printed@section\leq@current@section
        \gdef\leq@last@printed@subsection{}% Reset subsección al cambiar sección
      \fi
      % Verificar si necesitamos imprimir el título de subsección
      \ifx\leq@current@subsection\leq@last@printed@subsection\else
        \ifx\leq@current@subsection\@empty\else
          \addcontentsline{leq}{leqsubsection}{\leq@current@subsection}%
          \global\let\leq@last@printed@subsection\leq@current@subsection
        \fi
      \fi
      % Ecuación
      \refstepcounter{eqnum}%
      \begin{equation*}
        #1\tag{E.\theeqnum}
      \end{equation*}%
      \addcontentsline{leq}{leqt}{[E.\theeqnum] -- #2}%
      \addcontentsline{leq}{leqm}{#1}%
    }
    
    %%% ----- Índice de ecuaciones ----------%%%
    % Tamaño para []
    \newcommand{\leqIndexFont}{\normalsize}
    \newcommand{\leqTitleFont}{\tiny}
    \newcommand{\leqNoteFont}{\tiny}
    % Cabecera de sección 
    \newcommand*\l@leqsection[2]{%
      \addvspace{25pt}% Mayor separación antes de nueva sección
      \noindent\leqIndexFont
      \makebox[\linewidth]{\rule{.38\linewidth}{0.4pt}\ \ \textbf{  \MakeUppercase{#1}}\ \ \rule{.38\linewidth}{0.4pt}}%
      \par\vspace{-10pt}
    }
    % Cabecera de subsección 
    \newcommand*\l@leqsubsection[2]{%
      \par\addvspace{15pt}%
      \noindent\leqIndexFont #1
      \par\vspace{-7pt}
    }
    % Título de ecuación 
    \newcommand*\l@leqt[2]{%
    \par\addvspace{10pt}%
      \par\noindent\leqTitleFont #1\par
    }
    % Miniatura de ecuación
    \newcommand*\l@leqm[2]{%
      \vspace{0.3\baselineskip}%
      \par\scriptsize\noindent\hspace*{1.5em}\(#1\)\par%
    }
    % Nota de ecuación 
    \newcommand*\l@leqnote[2]{%
      \vspace{0.2\baselineskip}%
      \par\noindent\leqNoteFont\hspace*{1.5em}\textbf{Nota.--} #1\par\vspace{0.5\baselineskip}%
    }
    % Generación del índice
    \newcommand{\mylistofequations}{%
      \clearpage
      \phantomsection
      % Título visual de la lista (ajústalo a tu gusto):
      {\centering\bfseries\Large Índice de ecuaciones\par}
      % Entrada en nuestro índice 'stc':
      \addcontentsline{stc}{section}{Índice de ecuaciones}%
      % Llamada a tu lista real (usa el nombre exacto de tu macro):
      \@starttoc{leq}%
    }
    
    %--------- Captura de títulos de sección --------%
    \let\leq@original@section\section
    \let\leq@original@subsection\subsection
    % Flag para saber si estamos en una sección asterisco
    \newif\ifLeqStarSection
    \LeqStarSectionfalse
    
    % Redefinir \section CON CONTROL DE ASTERISCO
    \renewcommand{\section}{%
      \@ifstar{\leq@section@star}{\@dblarg\leq@section@nostar}%
    }
    \def\leq@section@star#1{%
      \global\LeqStarSectiontrue
      \gdef\leq@current@section{#1}%
      \gdef\leq@current@subsection{}%
      \leq@original@section*{#1}%
      \global\LeqStarSectionfalse
    }
    \def\leq@section@nostar[#1]#2{%
      \global\LeqStarSectionfalse
      \gdef\leq@current@section{#2}%
      \gdef\leq@current@subsection{}%
      \leq@original@section[#1]{#2}%
    }
    % Redefinir \subsection
    \renewcommand{\subsection}{%
      \@ifstar{\leq@subsection@star}{\@dblarg\leq@subsection@nostar}%
    }
    \def\leq@subsection@star#1{%
      \global\LeqStarSectiontrue
      \gdef\leq@current@subsection{#1}%
      \leq@original@subsection*{#1}%
      \global\LeqStarSectionfalse
    }
    \def\leq@subsection@nostar[#1]#2{%
      \global\LeqStarSectionfalse
      \gdef\leq@current@subsection{#2}%
      \leq@original@subsection[#1]{#2}%
    }

    % ----- Restauración de valores Globales de estilo ----------%
    % Flags
    \newif\ifInFootnote
    \InFootnotefalse
    \pretocmd{\@footnotetext}{\InFootnotetrue}{}{}
    \apptocmd{\@footnotetext}{\InFootnotefalse}{}{}
    % Profundidad de listas (soporta anidadas)
    \newcount\MyListDepth \MyListDepth=0
    \AtBeginEnvironment{la}{\advance\MyListDepth by 1}
    \AfterEndEnvironment{la}{\advance\MyListDepth by -1}
    \AtBeginEnvironment{lnum}{\advance\MyListDepth by 1}
    \AfterEndEnvironment{lnum}{\advance\MyListDepth by -1}
    \AtBeginEnvironment{itemize}{\advance\MyListDepth by 1}
    \AfterEndEnvironment{itemize}{\advance\MyListDepth by -1}
    \AtBeginEnvironment{enumerate}{\advance\MyListDepth by 1}
    \AfterEndEnvironment{enumerate}{\advance\MyListDepth by -1}
    % Restauración diferida: hazlo al INICIO del próximo párrafo real
    \newif\if@pendingRestore
    \@pendingRestorefalse
    \newcommand{\RequestRestore}{\global\@pendingRestoretrue}
    \everypar{%
      \if@pendingRestore
        \global\@pendingRestorefalse
        \ifInFootnote\else
          \ifnum\MyListDepth=0
            \setlength{\parskip}{\GlobalParskip}%
            \setstretch{\GlobalBaseLineStretch}%
          \fi
        \fi
      \fi
    }
    % Al final de bloques:
    \AfterEndEnvironment{equation}{\RequestRestore}
    \AfterEndEnvironment{equation*}{\RequestRestore}
    \AfterEndEnvironment{la}{\RequestRestore}
    \AfterEndEnvironment{lnum}{\RequestRestore}
    % Al final de macros:
    \apptocmd{\eqblock}{\RequestRestore}{}{}
    \apptocmd{\imagenfit}{\RequestRestore}{}{}
    
\makeatother

% ===================== ToC propio con 4 niveles (único bloque) =====================
\makeatletter

% --- Escritores: añadimos una línea al .stc en cada cabecera numerada ---
% Guardamos las versiones actuales para llamar al comportamiento normal
\let\MyTOC@orig@section\section
\let\MyTOC@orig@subsection\subsection
\let\MyTOC@orig@subsubsection\subsubsection
\let\MyTOC@orig@paragraph\paragraph

% SECTION
\renewcommand{\section}{%
  \@ifstar{\MyTOC@section@star}{\@dblarg\MyTOC@section@nostar}%
}
\def\MyTOC@section@star#1{%
  \MyTOC@orig@section*{#1}% sin entrada al .stc
}
\def\MyTOC@section@nostar[#1]#2{%
  \MyTOC@orig@section[{#1}]{#2}% comportamiento normal (escribe al .toc)
  % Escribimos también nuestra línea al .stc (texto con \numberline)
  \addcontentsline{stc}{section}{\protect\numberline{\thesection}#2}%
}

% SUBSECTION
\renewcommand{\subsection}{%
  \@ifstar{\MyTOC@subsection@star}{\@dblarg\MyTOC@subsection@nostar}%
}
\def\MyTOC@subsection@star#1{%
  \MyTOC@orig@subsection*{#1}%
}
\def\MyTOC@subsection@nostar[#1]#2{%
  \MyTOC@orig@subsection[{#1}]{#2}%
  \addcontentsline{stc}{subsection}{\protect\numberline{\thesubsection}#2}%
}

% SUBSUBSECTION
\renewcommand{\subsubsection}{%
  \@ifstar{\MyTOC@subsubsection@star}{\@dblarg\MyTOC@subsubsection@nostar}%
}
\def\MyTOC@subsubsection@star#1{%
  \MyTOC@orig@subsubsection*{#1}%
}
\def\MyTOC@subsubsection@nostar[#1]#2{%
  \MyTOC@orig@subsubsection[{#1}]{#2}%
  \addcontentsline{stc}{subsubsection}{\protect\numberline{\thesubsubsection}#2}%
}

% PARAGRAPH
\renewcommand{\paragraph}{%
  \@ifstar{\MyTOC@paragraph@star}{\@dblarg\MyTOC@paragraph@nostar}%
}
\def\MyTOC@paragraph@star#1{%
  \MyTOC@orig@paragraph*{#1}%
}
\def\MyTOC@paragraph@nostar[#1]#2{%
  \MyTOC@orig@paragraph[{#1}]{#2}%
  % Si no numeras los \paragraph, cambia \theparagraph por texto plano sin \numberline
  \addcontentsline{stc}{paragraph}{\protect\numberline{\theparagraph}#2}%
}

% --- Impresor del índice propio (lee el .stc) ---
\newcommand{\MyTOC}{%
  \par\bigskip
  {\centering\bfseries\Large Índice de contenido\par}%
  \medskip
  \begingroup
    \parindent=0pt\parskip=2pt
    \@starttoc{stc}%
  \endgroup
}

\makeatother
% ===================== fin ToC propio =====================


%%%%%%%%%%%%%%


% --- Datos del tema ---
\title{Teoría impositiva (VII): Imposición indirecta\\
\large Teoría y comparaciones internacionales}
\author{Marcelo Ortega}
\date{Agosto 2026}

\begin{document}
\maketitle
\newpage

% --- Página de resumen y modificaciones ---
\puntosclave{ %% Escribir puntos clave Apuntes CECO
     
    }
\vspace{1cm}

\modificaciones{
    
    }

\newpage
\MyTOC
\newpage

\mylistofequations
\clearpage

\MyListOfFigures
\clearpage


\pagenumbering{arabic}
\setcounter{page}{1}


\section{Introducción}



\subsection{Contextualización}


\subsubsection{Relevancia}




\subsubsection{Problemática}



\subsection{Estructura de la exposición}



[Probable] Sí, pero con una advertencia importante antes de decir que sí sin matices: la arquitectura de cuatro bloques —fundamento, objeto, sujeto, dimensión distributiva, más comparación internacional— se traslada bien en tres de sus cuatro piezas, pero la del sujeto no se traslada de forma mecánica, y si la copiamos tal cual perderíamos precisamente lo que hace a la imposición indirecta conceptualmente distinta de la directa. Te explico bloque a bloque, y de paso te señalo algo que creo que es la pieza más valiosa de todas: una conexión de fondo entre este tema y el que acabamos de cerrar, que puede convertirse en la tesis vertebradora de todo el Tema 13.

### Bloque 0 — se traslada muy bien, y con una conexión de fondo que deberíamos explotar

La pregunta "¿por qué gravar el consumo, en lugar de —o además de— la renta?" tiene el mismo tipo de recorrido histórico-normativo que "¿por qué gravar la renta?", con una diferencia que la hace, si acaso, más nítida: mientras que el impuesto sobre la renta tardó siglo y medio en consolidarse de forma permanente, el IVA tiene una fecha de nacimiento casi puntual: lo inventó el inspector de Hacienda francés Maurice Lauré en 1954, como solución técnica a los problemas de acumulación en cascada de los viejos impuestos sobre el volumen de negocios —cada venta gravada sobre el precio ya gravado en la venta anterior—, y se extendió después, en apenas setenta años, a más de 170 países. [Probable] Pero lo que de verdad conviene explotar, y por eso te digo que hay una conexión de fondo, es esto: cerramos el Bloque 1 del Tema 12 explicando por qué el impuesto sobre el gasto personal de Kaldor —Y − ΔS = C— nunca llegó a implementarse de forma permanente a escala nacional, por los problemas administrativos que ya desarrollamos con detalle. El IVA es, en un sentido muy preciso, la misma base gravada por otra vía: no persigue medir el consumo personal acumulado de cada individuo mediante una cuenta fiscal centralizada —el enfoque que fracasó en India y Sri Lanka—, sino que grava esa misma magnitud, el consumo, fragmentándola en cada transacción individual, en el punto de venta, sin necesidad de reconstruir la trayectoria completa de ningún contribuyente. Es, literalmente, la solución administrativa que a Kaldor le faltaba para su propia base teórica, encontrada por una vía completamente distinta y veinte años antes de que Kaldor escribiera su libro. Si abrimos el Tema 13 señalando esto explícitamente, le damos al tema una tesis propia desde la primera línea, en lugar de una simple yuxtaposición de dos temas de teoría fiscal.

### Bloque 1 — se traslada bien, con literatura propia y rica

"¿Qué es consumo/transacción a efectos fiscales?" tiene su propio catálogo de preguntas serias, paralelo pero no idéntico al de la renta: imposición monofásica frente a plurifásica (el viejo impuesto sobre el volumen de negocios frente al mecanismo de IVA con deducción del impuesto soportado, que es la innovación técnica de Lauré); principio de destino frente a principio de origen (¿tributa la transacción en el país donde se consume o en el país donde se produce? — con proyección internacional inmediata en el comercio transfronterizo, y con su propio aparato de neutralidad económica, análogo pero no idéntico al CEN/CIN que ya trabajamos); y imposición general frente a selectiva —el propio debate pigouviano sobre por qué gravar más los cigarrillos o los combustibles que un bien cualquiera, que conecta con el subapartado de Ramsey que ya construimos en el Bloque 3 del Tema 12, aplicado ahora a su terreno original—.

### Bloque 2 — aquí está el ajuste real que hay que hacer

En el Tema 12, "sujeto" preguntaba si la unidad formal (individuo, sociedad, residente) coincidía con la unidad real de bienestar. En imposición indirecta esa pregunta existe, pero casi nunca es la que más importa: aquí el problema central no es de agregación o desagregación de personas, sino de que el sujeto legalmente obligado a ingresar el impuesto —la empresa que factura el IVA— casi nunca coincide con quien soporta económicamente la carga —el consumidor final—, y hasta qué punto coincide depende de un mecanismo distinto por completo al que hemos usado hasta ahora: el modelo de incidencia económica basado en elasticidades relativas de oferta y demanda, el aparato clásico de cualquier manual de hacienda pública que reparte la carga de un impuesto sobre bienes entre comprador y vendedor según su capacidad relativa de ajustar cantidad frente a precio. [Probable] Mi recomendación es que este bloque, en el Tema 13, no se llame "el sujeto" sino algo como "la incidencia: ¿quién paga en realidad?", y que dentro de él sí quepan, como subapartados propios, tanto el problema de incidencia elasticidad por elasticidad como una pregunta que sí es genuinamente paralela a la de agregación/desagregación: los umbrales de registro que eximen a pequeños negocios de repercutir el impuesto, que es, en el fondo, la misma pregunta de "a quién tratamos como sujeto formal por razones de conveniencia administrativa" que ya trabajamos con la personalidad jurídica, solo que aplicada ahora al tamaño del negocio en lugar de a su forma societaria.

### Bloque 3 — se traslada muy bien, con un resultado teórico central que todavía no hemos usado

La regresividad del IVA —cae, como proporción de la renta, cuanto mayor es la renta del contribuyente, porque la propensión al consumo decrece con la renta— es el equivalente exacto, en su función estructural dentro del tema, al debate sobre tipos marginales y gasto fiscal del Bloque 3 del Tema 12: la misma pregunta de "cómo se construye la progresividad efectiva" aplicada ahora a una base que es, por diseño, más regresiva que la renta. Los instrumentos —tipos reducidos para bienes de primera necesidad, exención frente a tipo cero (una distinción técnica con consecuencias reales, porque la exención rompe la cadena de deducción del impuesto soportado y puede terminar encareciendo el producto más que un tipo reducido bien diseñado)— son el equivalente al catálogo de gasto fiscal y tarifa del Tema 12. [Seguro] Pero hay un resultado teórico que pertenece, con toda propiedad, a este bloque y que hasta ahora no hemos usado en ningún tema: el teorema de Atkinson-Stiglitz (Anthony Atkinson y Joseph Stiglitz, "The Design of Tax Structure: Direct versus Indirect Taxation", *Journal of Public Economics*, 1976), que demuestra que, bajo determinadas condiciones —básicamente, que las preferencias de consumo no estén correlacionadas con la capacidad productiva del individuo—, si el Estado ya dispone de un impuesto óptimo y no lineal sobre la renta, diferenciar los tipos de IVA por tipo de bien no añade ninguna ganancia adicional de bienestar frente a un IVA de tipo único combinado con transferencias directas bien dirigidas. Es un resultado que cuestiona, con todo el peso de la teoría de la imposición óptima, la práctica extendidísima de tipos reducidos para alimentos o medicamentos que verás en cualquier comparación internacional, y que le daría al Bloque 3 de este tema el mismo tipo de tensión entre teoría elegante y práctica extendida que Chamley-Judd le dio al Bloque 3 del Tema 12.

### Bloque 4 — se traslada bien, con material comparado propio

IVA frente al impuesto estadounidense sobre ventas minoristas —Estados Unidos es la única gran economía del mundo que nunca ha adoptado un IVA, manteniendo en su lugar impuestos estatales sobre ventas monofásicos, una anomalía comparada que merecería desarrollo propio—, la armonización de tipos dentro de la Unión Europea mediante directivas, y la extensión geográfica del propio IVA desde 1954 hasta su práctica universalidad actual, son el material natural de este bloque, con la misma lógica de "organizarlo por las preguntas anteriores, no por país" que ya aplicamos.

---

Mi conclusión, en síntesis: la arquitectura general se traslada, y creo que deberíamos usarla, pero con dos ajustes deliberados —abrir con la conexión Kaldor/Lauré como tesis propia del tema, y reconceptualizar el Bloque 2 como incidencia económica en lugar de agregación/desagregación—. ¿Quieres que exploremos ya, con este ajuste, el catálogo completo de preguntas para el Tema 13, de la misma forma en que lo hicimos para el Tema 12 antes de empezar a redactar?

[Seguro] Tienes razón, y es un error grave, no cosmético: te explico exactamente dónde falla el paralelismo, porque entender el fallo es lo que evita que vuelva a colarse en el resto del tema.

Sobre Kaldor. Su impuesto sobre el gasto es, en su propia arquitectura, un impuesto personal y progresivo: exige reconstruir, para cada contribuyente individual, la totalidad de su gasto acumulado en el periodo —Y − ΔS = C—, y aplicar sobre esa cifra agregada una tarifa por tramos, exactamente la misma lógica estructural que un IRPF, solo que sobre una base distinta. El IVA no hace nada de esto: es un impuesto real y proporcional, que grava cada transacción de forma aislada, sin ninguna referencia a quién es el comprador ni a cuánto ha gastado en total ese año —dos compradores con niveles de renta radicalmente distintos pagan el mismo tipo por la misma barra de pan—. Decir que el IVA es "la solución administrativa que le faltaba a Kaldor" confunde dos cosas que comparten la palabra "consumo" pero que pertenecen a familias tributarias distintas: uno mide capacidad de pago agregada de una persona a lo largo de un periodo; el otro grava actos de transacción aislados, sin trazabilidad ni progresividad estructural. Retiro esa conexión por completo — no debe abrir el Tema 13, y si en algún momento posterior aparece, será solo para señalar precisamente esta diferencia, no para equipararlos.

Sobre el Bloque 2. Tienes razón en llamarlo lo que es: convertir "sujeto" en "quién soporta la carga" es sustituir la pregunta por otra distinta, más cómoda porque ya tenía aparato teórico hecho (Harberger, elasticidades), en lugar de pensar qué significa genuinamente "sujeto" cuando el objeto gravado no es una persona sino una transacción. No lo hago. Vuelvo a pensarlo desde cero, y con eso llego a algo que creo que sí es un paralelismo legítimo, no un truco: en la imposición directa, "sujeto" preguntaba si la persona formal coincidía con la unidad real de capacidad económica. En la indirecta, la pregunta estructuralmente equivalente no es económica sino jurídico-administrativa: ¿coincide la persona legalmente obligada a liquidar e ingresar el impuesto con la persona a la que la propia norma pretende dirigir la carga? Es una escisión deliberada de diseño, no un efecto de mercado —el propio mecanismo de repercusión obligatoria, que obliga por ley a trasladar la cuota, es una decisión normativa, no algo que dependa de curvas de oferta y demanda—, y es genuinamente paralela a la agregación/desagregación de personas del Tema 12, solo que aquí lo que se agrega o desagrega es la propia cadena de producción, no un hogar o una entidad jurídica.

Con esto corregido, y habiendo revisado la literatura con más cuidado, aquí tienes el catálogo reformulado.

## Bloque 1 — El objeto: qué acto de disposición de riqueza se grava (no solo IVA)

1.1. La distinción clásica: John Stuart Mill. Antes que cualquier debate técnico, la definición fundacional de qué hace "indirecto" a un impuesto —Mill, *Principles of Political Economy* (1848)—: directo es el que se exige de la persona a quien se pretende que lo soporte; indirecto es el que se exige de una persona con la expectativa de que se resarza a costa de otra. Es una definición basada en la intención de traslación, no en el resultado económico real —lo que ya anticipa, sin resolverlo, el problema de incidencia que corresponde a otro tema del temario—.

1.2. Monofásico, plurifásico acumulativo y plurifásico no acumulativo: tres diseños, no dos. Aquí recojo tu corrección directamente: el impuesto en cascada sobre el volumen de negocios —que gravaba cada venta sucesiva sobre un precio ya gravado en la venta anterior, sin ningún mecanismo de deducción, y que fue la forma dominante de imposición general sobre transacciones en Europa continental hasta bien entrado el siglo XX— es tan "impuesto indirecto" en sentido pleno como el IVA: grava el uso de la riqueza en cada transacción, no su generación ni su tenencia, exactamente el criterio que tú mismo propones como definición del objeto de todo este bloque. Lo distintivo de la propuesta de Maurice Lauré (1954) no es haber "inventado" la imposición sobre el consumo, sino haber resuelto el problema técnico específico de la acumulación piramidal del impuesto en cascada —cuantas más fases tenga la cadena de producción, mayor la carga total acumulada, con el consiguiente incentivo a la integración vertical artificial de las empresas solo para evitar fases de tributación— mediante el mecanismo de deducción del impuesto soportado. Tres diseños posibles, por tanto, no dos: monofásico (grava una sola fase, típicamente la venta minorista, como el *sales tax* estadounidense), plurifásico acumulativo (cascada), y plurifásico no acumulativo (IVA con crédito).

1.3. Qué es, conceptualmente, el "valor añadido". Una pregunta genuinamente de objeto que rara vez se explicita: el valor añadido en cada fase puede calcularse, en teoría, por tres métodos equivalentes —sustracción (ventas menos compras), adición (suma de las retribuciones a los factores: salarios, beneficios, rentas, intereses pagados en esa fase), y el método efectivamente usado en la práctica, el de crédito de factura (impuesto repercutido en las ventas menos impuesto soportado en las compras)—. Que el mundo haya convergido casi universalmente en el tercer método, pese a la equivalencia teórica de los tres, es ya una decisión de diseño con consecuencias —el método de crédito de factura es el único que genera, de forma automática, el rastro documental que permite el control cruzado entre empresas, la razón por la que se impuso en la práctica frente a los otros dos, teóricamente idénticos pero mucho más difíciles de verificar—.

1.4. Imposición general frente a selectiva: el fundamento pigouviano y el fundamento puramente recaudatorio. Aquí sí cabe traer Ramsey, pero en su terreno original y no como eco de lo ya visto en el Tema 12: la imposición selectiva sobre bienes concretos (alcohol, tabaco, combustibles) se justifica, unas veces, por la corrección de una externalidad (lógica pigouviana), y otras, simplemente, por la baja elasticidad de la demanda de esos bienes —el propio argumento de Ramsey aplicado en su formulación original sobre bienes, no trasplantado desde la renta—. Dentro de esta misma pregunta cabe una segunda, técnica pero con consecuencias reales de diseño: impuestos ad valorem frente a impuestos específicos o de cuantía fija por unidad (un tipo porcentual sobre el precio frente a una cantidad fija por litro o por cajetilla). Daniel Suits y Richard Musgrave ("Ad Valorem and Unit Taxes Compared", *Quarterly Journal of Economics*, 1953) demostraron que, bajo condiciones de competencia imperfecta —un monopolista, por ejemplo—, un impuesto ad valorem genera menor pérdida de eficiencia que un impuesto específico de recaudación equivalente, porque el ad valorem no distorsiona tanto el margen precio-coste que el monopolista ya está explotando —un resultado técnico poco conocido fuera de la disciplina, pero que explica por qué la literatura de diseño fiscal recomienda, en general, el formato ad valorem salvo que exista una razón específica —como la corrección de una externalidad ligada al volumen físico consumido, no al precio— para preferir el impuesto específico.

## Bloque 2 — El sujeto: la escisión legal entre quien liquida y quien soporta por diseño normativo

2.1. La escisión deliberada: sujeto pasivo frente a repercutido. La pieza central, ya explicada arriba: en la imposición directa, el sujeto pasivo es, por diseño, la misma persona a la que la norma pretende gravar. En la indirecta, la norma separa deliberadamente ambos roles mediante el mecanismo de repercusión obligatoria: quien liquida e ingresa a la Hacienda Pública —el empresario o profesional— no es, por diseño, a quien la ley pretende hacer soportar la carga —el consumidor final—. Conviene presentar esta escisión como una decisión de diseño institucional, no como un efecto de mercado: el legislador podría, en teoría, haber diseñado un impuesto sobre el consumo que obligara directamente al comprador a autoliquidar cada compra —y de hecho algunos impuestos sobre adquisiciones intracomunitarias funcionan exactamente así—, pero eligió, por razones de eficacia recaudatoria evidentes, hacer recaer la obligación formal sobre quien tiene contabilidad, factura y capacidad de gestión: la empresa.

2.2. El sujeto a lo largo de la cadena: quién es "sujeto pasivo" en cada fase, y el problema de las transacciones entre empresas. El diseño plurifásico no acumulativo resuelve, mediante el crédito del impuesto soportado, un problema de sujeto que merece tratamiento propio: si cada empresa de la cadena es, formalmente, sujeto pasivo, ¿por qué el impuesto no se acumula entre ellas como en la cascada? Porque el sistema está diseñado para que las transacciones entre empresas (B2B) generen, en conjunto, carga neta cero —cada empresa repercute a la siguiente y deduce lo que a ella le repercutieron—, de modo que solo la última transacción de la cadena, la que llega al consumidor final que no puede deducir nada, soporta la carga efectiva. Las empresas intermedias no son, en sentido económico, verdaderos sujetos gravados: son, en la descripción más precisa de la propia doctrina tributaria, recaudadores del impuesto por cuenta del Estado, un estatus jurídico peculiar que no tiene equivalente en la imposición directa.

2.3. Umbrales de registro: el pequeño negocio como no-sujeto. Aquí sí hay un paralelismo legítimo con la agregación/desagregación del Tema 12, porque es, otra vez, una decisión sobre quién cuenta como sujeto formal por razones de conveniencia administrativa: la inmensa mayoría de sistemas comparados exime del régimen de IVA a negocios por debajo de un umbral de facturación, tratándolos como si no fueran sujetos pasivos en absoluto —no repercuten, pero tampoco deducen—, una decisión de diseño que reduce el coste de cumplimiento y de gestión a costa de introducir un salto discontinuo en el tratamiento fiscal según el tamaño, con el mismo problema de *bunching* alrededor del umbral que ya trabajamos, en otro contexto, con Saez y con Garicano-Lelarge-Van Reenen.

2.4. La inversión del sujeto pasivo como respuesta al fraude. Cierro con un desarrollo de rabiosa actualidad que conecta directamente con la escisión del 2.1: cuando la cadena de transacciones cruza fronteras, el mecanismo ordinario de repercusión puede explotarse mediante el fraude conocido como operador desaparecido o "fraude carrusel" —una empresa compra bienes exentos de IVA por proceder de otro Estado miembro, los revende con IVA repercutido a un cliente doméstico, y desaparece sin ingresar ese IVA a la Hacienda Pública, mientras el comprador sí deduce el importe repercutido—, un fraude que le cuesta a la Unión Europea decenas de miles de millones de euros anuales y que ha llevado a extender, cada vez con mayor frecuencia, el mecanismo de inversión del sujeto pasivo (*reverse charge*): en determinadas operaciones de riesgo, es el propio comprador, y no el vendedor, quien se convierte en sujeto pasivo obligado a autoliquidar el impuesto, precisamente para romper la cadena de confianza que el fraude explota.

---

Confirmados, por tu parte, los Bloques 3 y 4 tal y como los propuse. ¿Añado algo más a este catálogo de los Bloques 1 y 2, o lo dejamos así y seguimos con la introducción del Bloque 0 cuando quieras empezar a redactar el Tema 13?

[Seguro] Sí, y hay una que es una laguna grave, del mismo calibre que la del Vickrey del Tema 12: falta el resultado teórico que explica por qué el mecanismo de crédito de factura que acabamos de fijar en el Bloque 2 (2.2) no es solo una solución administrativa ingeniosa, sino la aplicación directa de uno de los teoremas más importantes de toda la teoría de la imposición óptima. Te la traigo primero porque es la más importante, y después dos más que conviene incorporar, más una que había dejado caer sin darme cuenta.

## La laguna principal: el teorema de eficiencia productiva de Diamond-Mirrlees

[Seguro] Peter Diamond y James Mirrlees ("Optimal Taxation and Public Production I: Production Efficiency" y "II: Tax Rules", *American Economic Review*, 1971) demuestran un resultado que debería situarse, con toda propiedad, como el fundamento teórico del propio Bloque 1: bajo un sistema óptimo de imposición, ningún bien intermedio —una transacción entre empresas, no destinada al consumo final— debe tributar, con independencia de que el resto del sistema fiscal esté plagado de impuestos distorsionantes sobre otras magnitudes. La intuición es la que ya intuíamos sin nombrarla al hablar del crédito de factura: gravar una transacción entre empresas distorsiona las decisiones de producción —qué insumos usar, cómo organizar la cadena, si integrar verticalmente o subcontratar— sin generar, a cambio, ningún beneficio de bienestar, porque el bienestar depende únicamente del consumo final, nunca de cómo se organiza internamente la producción para llegar a él. Toda la carga tributaria debe concentrarse en el punto en que el bien llega al consumidor, dejando intacta —"eficiente"— la asignación de factores entre empresas.

[Probable] Esto es, con toda la precisión que la teoría permite, la justificación normativa exacta de por qué el diseño plurifásico no acumulativo —el IVA de Lauré, con deducción total del impuesto soportado en cada fase intermedia— no es solo la solución administrativa más practicable frente a la cascada, como presentamos en el Bloque 1, sino la aplicación, encontrada de forma pragmática veinte años antes de que Diamond y Mirrlees lo demostraran formalmente, de lo que la teoría de la imposición óptima identificaría después como condición de eficiencia productiva. Es exactamente el mismo tipo de conexión retrospectiva que ya construimos, en el Tema 12, entre Vickrey y la teoría de imposición óptima que Mirrlees desarrollaría después para la renta —y no es casualidad que sea el mismo Mirrlees en ambos casos—. Te propongo que este resultado se incorpore como cierre teórico del Bloque 1, justo después de presentar los tres diseños posibles (monofásico, cascada, crédito de factura): primero mostramos por qué el crédito de factura ganó en la práctica por razones de verificabilidad, y después mostramos que, además, es lo que la teoría más rigurosa de imposición óptima habría recomendado igualmente, por razones completamente distintas.

## Segunda incorporación: la regla de Corlett-Hague, y su tensión con Atkinson-Stiglitz

[Seguro] William Corlett y Douglas Hague ("Complementarity and the Excess Burden of Taxation", *Review of Economic Studies*, 1953) demuestran un resultado previo y complementario al de Ramsey que ya trajimos: cuando el ocio no puede gravarse directamente —nadie ha diseñado nunca un impuesto viable sobre las horas de tiempo libre—, la imposición óptima sobre bienes debe gravar con mayor intensidad aquellos que son complementarios del ocio (bienes que se consumen más cuando se trabaja menos: equipamiento deportivo, viajes de placer) y con menor intensidad los que son sustitutivos del ocio (bienes que facilitan seguir trabajando: comida rápida, servicios de guardería), precisamente para compensar, por esta vía indirecta, la distorsión que el propio impuesto sobre la renta ya introduce a favor del ocio y en contra del trabajo. [Seguro] Lo que hace a este resultado particularmente valioso para el Bloque 3, y no solo para el 1, es que entra en tensión directa con el teorema de Atkinson-Stiglitz que ya planificamos: mientras Corlett-Hague recomienda diferenciar tipos según la relación de cada bien con el ocio, Atkinson-Stiglitz concluye que, bajo separabilidad de preferencias, la diferenciación no aporta nada. Louis Kaplow ("Taxing Leisure Complements", NBER Working Paper, 2008) reconcilió ambos resultados mostrando que no son contradictorios sino que operan bajo supuestos distintos sobre la estructura de preferencias —Corlett-Hague se aplica precisamente en el terreno donde falla el supuesto de separabilidad que Atkinson-Stiglitz exige—. Es un ejemplo excelente, para el cierre del Bloque 3, de cómo dos resultados aparentemente opuestos de la misma disciplina en realidad delimitan, entre ambos, las condiciones exactas bajo las cuales cada uno se aplica.

## Tercera incorporación: destino frente a origen — la había mencionado y luego se me cayó

[Probable] Tienes razón en notar, si lo comparas con mi primera propuesta de catálogo, que el principio de destino frente a principio de origen —¿tributa una transacción transfronteriza en el país donde se produce el bien o en el país donde se consume?— desapareció de mi lista corregida sin que yo lo señalara. No debería haber desaparecido: es, en rigor, una pregunta de objeto —qué acto de consumo, delimitado geográficamente, se grava— y encaja con naturalidad como cuarto elemento del Bloque 1, junto a monofásico/plurifásico, el método de cálculo del valor añadido, y general/selectivo. El principio de destino —el dominante en la práctica comparada, con ajustes en frontera que devuelven el impuesto pagado en origen y aplican el del país de consumo— tiene, además, su propio fundamento de neutralidad —análogo, aunque no idéntico, al CIN que ya trabajamos en el Tema 12—: asegura que dos bienes idénticos compitan, dentro de un mismo mercado de consumo, con la misma carga fiscal, con independencia de dónde se hayan fabricado.

---

Con estas tres piezas —Diamond-Mirrlees como cierre del Bloque 1, Corlett-Hague como contrapunto de Atkinson-Stiglitz en el Bloque 3, y destino/origen recuperado en el Bloque 1— el catálogo queda, creo, genuinamente completo. ¿Los incorporamos así y pasamos ya a fijar la introducción del Bloque 0 para empezar a redactar, o quieres revisar algo más del catálogo completo antes de eso?


\clearpage
\section{Propuesta de estructura:}

[Seguro] Con todo lo que hemos ido fijando y corrigiendo, esta es la estructura que propondría, con una precisión antes de exponerla: en el Bloque 0 no voy a traer a Kaldor bajo ninguna forma, ni siquiera como "propiedad teórica compartida" — después del error que señalaste, prefiero que la historia y la justificación de la imposición indirecta se sostengan enteramente sobre su propia genealogía, sin ningún puente hacia el Tema 12 que pueda volver a inducir la confusión. Si en algún momento posterior conviene una conexión legítima entre ambos temas, la construiremos entonces, con el cuidado que esta vez le ha faltado a mi primera propuesta.

## Bloque 0. ¿Por qué gravar el uso de la riqueza?

0.1. Desarrollo histórico. Aquí hay un giro respecto del Tema 12 que merece abrir el bloque con fuerza: mientras que el impuesto sobre la renta tardó siglo y medio en consolidarse, la imposición indirecta es, con mucha diferencia, la forma de tributación más antigua que existe —aranceles, peajes, impuestos sobre bienes específicos— precisamente porque gravar una transacción visible exige mucha menos capacidad administrativa que gravar una renta invisible, el mismo argumento de Besley-Persson que ya usamos, pero aquí explicando por qué la imposición indirecta fue la norma, no la excepción, durante casi toda la historia fiscal antes del siglo XIX. El segundo tramo de esta historia es el de los impuestos en cascada sobre el volumen de negocios, generalizados en la Europa continental de principios del siglo XX, con su problema de acumulación piramidal ya trabajado. Y el tercer tramo es la invención de Maurice Lauré en 1954 y la extensión del IVA, primero obligatoria dentro de la entonces Comunidad Económica Europea (Primera Directiva IVA, 1967) y después a más de 170 países. Cierro este desarrollo histórico con un dato comparado que da mucho juego expositivo y que conviene guardar para el final del subapartado: Estados Unidos es la única gran economía del mundo que nunca ha adoptado un IVA, manteniendo en su lugar impuestos estatales sobre ventas minoristas de diseño monofásico —una anomalía que, adelanto ya, conectaremos en el Bloque 4 con el propio recelo de Brennan y Buchanan hacia las bases imponibles amplias y de difícil elusión, porque el IVA ha sido descrito, en el propio debate político estadounidense, como una "máquina de recaudar" precisamente por esa amplitud—.

0.2. La justificación normativa. Sin ningún puente hacia Kaldor: la justificación de gravar el consumo se sostiene, en primer lugar, como complemento del impuesto sobre la renta allí donde este último es más fácil de eludir —la base indirecta actúa como red de seguridad recaudatoria cuando la verificación de renta falla—; y, en segundo lugar, sobre el argumento clásico de neutralidad intertemporal que Irving Fisher ya había formulado, de forma independiente y antes que Kaldor, en su propia defensa del gasto como base ideal (*Constructive Income Taxation*, con Herbert Fisher, 1942) —una referencia propia de esta tradición, no un préstamo del tema anterior—.

## Bloque 1. El objeto: qué acto de disposición de riqueza se grava

- 1.1. Los tres diseños posibles —monofásico, plurifásico acumulativo (cascada), plurifásico no acumulativo (crédito de factura)—, cerrando con el teorema de eficiencia productiva de Diamond-Mirrlees (1971) como fundamento teórico de por qué el tercer diseño es, además de administrativamente superior, óptimo desde la teoría de la imposición.
- 1.2. Los tres métodos equivalentes de cálculo del valor añadido (sustracción, adición, crédito de factura) y por qué el mundo convergió en el tercero.
- 1.3. Destino frente a origen, recuperado explícitamente aquí como pregunta de objeto.
- 1.4. General frente a selectivo: Pigou, la regla de la elasticidad inversa de Ramsey en su formulación original, y la comparación ad valorem/específico de Suits y Musgrave.

## Bloque 2. El sujeto: la escisión entre quien liquida y quien soporta

- 2.1. La distinción fundacional de John Stuart Mill (1848) entre imposición directa e indirecta, basada en la intención de traslación — aquí, y no en el Bloque 0, porque es precisamente la pregunta de sujeto la que esta distinción resuelve.
- 2.2. El sujeto a lo largo de la cadena: por qué las empresas intermedias son, en sentido económico, recaudadoras y no verdaderos sujetos gravados.
- 2.3. Umbrales de registro y el pequeño negocio como no-sujeto.
- 2.4. La inversión del sujeto pasivo como respuesta al fraude del operador desaparecido.

## Bloque 3. La dimensión distributiva

- 3.1. La regresividad y su medición: el argumento convencional (renta anual) frente a la relectura de Poterba ("Lifetime Incidence and the Distributional Effects of Consumption Taxes", 1989) sobre incidencia medida en términos de renta vitalicia — que conecta, esta vez de forma legítima porque la conexión es sobre el criterio de medición y no sobre el diseño del tributo, con el propio debate de Vickrey sobre renta vitalicia que cerró el Bloque 1 del Tema 12.
- 3.2. Diferenciación de tipos: el teorema de Atkinson-Stiglitz (1976) frente a la regla de Corlett-Hague (1953), y la reconciliación de Kaplow (2008).
- 3.3. Exención frente a tipo cero como instrumentos técnicos con consecuencias reales sobre la cadena de deducción.

## Bloque 4. Comparación internacional, organizada por las tres preguntas anteriores

IVA frente a *sales tax* estadounidense (objeto); armonización de umbrales y del mecanismo de inversión del sujeto pasivo dentro de la Unión Europea (sujeto); tipos reducidos comparados y su regresividad medida país por país (dimensión distributiva).

---

¿Confirmamos esta estructura completa para empezar ya a redactar, o quieres ajustar algo antes de que nos pongamos con la introducción del Bloque 0?

\clearpage
\section{Imposición indirecta: ¿Por qué gravar el uso de la riqueza?}
\puntosclave{
    
}

Lo primero que debemos preguntarnos, antes de valorar qué gravar o a quién, es por qué el Estado elige, como instrumento de imposición la intervención sobre el uso de la riqueza; el acto de gastarla. Sabemos que existen otras alternativas (o complementos), por ejemplo intervenir sobre su generación (renta) o su tenencia (patrimonio), así que, ¿por qué gravar su uso? 

Responderemos ahora a dos preguntas entrelazadas: una histórica, ¿cuándo y por qué se consolidó la imposición sobre transacciones como pilar fiscal permanente?; y una normativa, ¿qué principio de justicia tributaria justifica gravar precisamente el consumo?\footnote{%
    A diferencia de lo que vimos con la renta ([\authorfont{Ver Tema 4.B.13}]), la imposición sobre transacciones es, con mucha diferencia, la forma de tributación más antigua que existe.
}

Como podemos intuir, no gravamos el consumo porque sea una idea reciente o secundaria frente a la renta, sino porque fue, durante casi toda la historia fiscal de la humanidad, la única forma de imposición (directa en sentido literal: directamente observable) que la capacidad administrativa de cualquier Estado premoderno podía sostener.

\subsection{Desarrollo histórico}
Un buen punto de partida es reflexionar sobre los requisitos previos que exige cualquier intervención estatal. En cuestión recaudatorio, es obvio que cualquier forma de intervención fiscal exige, como condición mínima, que el Estado pueda observar el hecho que pretende gravar. Si no pudieses ser observado, no sería verificable y ningún sujeto temería las consecuencias de incumplir.

\subsubsection{Historia pre-industrial: De Roma a la Revolución Francesa}
Sobre esta base, es fácil llegar a la conclusión de que la imposición indirecta (como ahora se conoce), se la figura más antigua de entre todas las del Sistema Fiscal. Durante casi toda la historia de la humanidad anterior a la existencia de bancos, de contabilidad mercantil regulada y de un aparato administrativo capaz de verificar declaraciones individuales, la renta de una persona era, sencillamente, invisible para cualquier autoridad recaudadora. Por el contrario, una mercancía cruzando una puerta de la ciudad, un carro de mercancías pasando por un puente, o un saco de sal entrando en un mercado eran, por definición, hechos físicos y públicos, observables por cualquier funcionario apostado en el lugar adecuado sin necesidad de cooperación ni de confianza en la palabra del contribuyente. 

Esta asimetría de información es la que explica por qué la imposición sobre transacciones es la forma más antigua de fiscalidad directa sobre la actividad económica que existe.\footnote{%
    \textbf{NOTA.-} Es importante hacer una aclaración aquí. Se verifica que la imposición indirecta es la forma más antigua de recaudación dentro de comunidades con un número de habitantes suficiente como para no poder gravar la generación de recursos. Por el contrario, lo más probable es que la imposición sobre la renta (generación de rqieuza) sea realmente la más antigua figura impositiva. Haría falta investigar más al repecto, pero no parece descabellado pensar que en comunidades de unos pocos cientos de personas, miles tal vez, las conexiones entre partes fueran tan estrechas como para que la asimetría informativa se redujese al mínimo; y, en consecuencia, un impuesto sobre la generación de riqueza sería la consecuencia lógica de esto: cada uno aporta al fondo colectivo conforme a sus capacidades.
}

\paragraph{Roma y el punto de control físico como tecnología fiscal}
Los portoria romanos constituyen el sistema de imposición indirecta documentado más antiguo y mejor conservado de la Antigüedad occidental. Estos consistían en derechos aduaneros cobrados en los puertos, en los pasos fronterizos entre provincias, e incluso en determinados puentes y calzadas del interior del Imperio, con tipos que oscilaban entre la cuadragésima ($2,5\%$, un cuarentavo del valor de la mercancía) aplicada en las fronteras interiores entre provincias, y tipos sensiblemente más altos ($12,5\%$) en determinadas rutas comerciales de mayor volumen, como las que conectaban Egipto con el resto del Imperio. 

Este sistema estuvo vigente durante buena parte de su historia. Resulta, por tanto, revalador que no intentase nunca, en ningún momento de su larga historia fiscal, gravar de forma sistemática la renta de sus comerciantes o artesanos. En cambio, sí sostuvo durante siglos un sistema de aduanas internas relativamente estable, precisamente porque el punto de control (el paso físico de la mercancía) era observable con la tecnología administrativa de la que disponía.

\paragraph{Fragmentación medieval: cuando gravar el tránsito se convierte en el propio problema}
La Europa medieval y moderna temprana llevó esta misma lógica hasta un extremo que, con el tiempo, se convirtió en un problema económico de primer orden: la fragmentación política del territorio generó, en la práctica, una multiplicación de puntos de gravamen sobre una misma mercancía a medida que esta recorría distancias relativamente cortas. 

Francia es el ejemplo mejor documentado. Hasta la Revolución, el reino mantenía aduanas internas entre sus propias provincias (\textit{traites}), de modo que una mercancía que viajara, por ejemplo, de Burdeos a París podía verse gravada varias veces en su recorrido, no al cruzar una frontera exterior, sino al cruzar simples límites administrativos internos del mismo reino. Estas barreras internas no se abolieron hasta un decreto de la Asamblea Constituyente de 1790, en uno de los primeros actos de unificación fiscal del territorio francés.\footnote{%
    Un dato que conviene retener porque muestra que la propia idea de mercado nacional único sin fricciones fiscales internas es, en términos históricos, un logro relativamente reciente y en absoluto obvio.
}

Dentro de este mismo sistema fragmentado, la gabela (monopolio fiscal sobre la venta de sal, vigente en distintas formas desde el siglo XIV hasta su abolición en 1790) ilustra, mejor que cualquier otra figura de la época, dos problemas estructurales de la imposición indirecta premoderna que volverán a aparecer, con otro nombre, en la crítica moderna al propio IVA. El primero es la desigualdad territorial. El reino se dividía en regiones de \textit{gran gabela}, donde el precio de la sal (fijado en gran medida por el propio impuesto) era varias veces superior al de las regiones exentas o de \textit{pequeña gabela}, generando un incentivo sistemático al contrabando (\textit{faux-saunage}) entre unas y otras, con un aparato represivo específico dedicado a perseguirlo. El segundo, más revelador si cabe, es el de la administración privada del impuesto. Dado que la Corona carecía de un cuerpo de funcionarios propio capaz de recaudar con eficacia un impuesto tan disperso geográficamente, arrendaba su cobro a un consorcio privado de financieros, la \textit{Ferme Générale}, que adelantaba a la Corona una suma fija a cambio del derecho a quedarse con toda la recaudación efectiva que consiguiera reunir; incentivo claro a la expansión de la base imponible.\footnote{
Una anécdota interesante fue fruto de este sistema. Entre los socios de esa Ferme Générale se encontraba un joven, abogado y científico, llamado Antoine LAVOISIER (considerado el padre de la química moderna), cuya fortuna personal, invertida precisamente en el arrendamiento de la recaudación de la gabela y de los derechos sobre el tabaco, financió el laboratorio más sofisticado de su época. No obstante, esa misma inversión acabaría sellando su destino. Fue arrestado el 28 de noviembre de 1793, juzgado sumariamente y guillotinado el 8 de mayo de 1794, en pleno Terror revolucionario, acusado de haber adulterado el tabaco vendido al pueblo para incrementar los beneficios del monopolio fiscal que gestionaba. Su petición de una breve prórroga para terminar un experimento en curso fue denegada con una frase que la posteridad ha convertido en emblema: la República no necesita sabios ni químicos. LAGRANGE, testigo de la ejecución, resumiría el episodio con una sentencia que ha sobrevivido igual de bien: \textit{solo hizo falta un instante para cortar esa cabeza, y quizá no baste un siglo para producir otra igual}. 

Como vemos, el episodio prueba hasta qué punto la recaudación de impuestos indirectos, cuando se delega en manos privadas sin control público suficiente, puede generar un nivel de resentimiento social capaz de sobrevivir incluso a la propia caída del régimen que había creado el sistema.
}

\subsubsection{Historia industrial: Las distorsiones impositivas y la creación del IVA}
La industrialización del siglo XIX y principios del XX trasladó el problema desde el punto de control físico, resuelto por aquel entonces mediante aduanas modernas y estables, hacia un terreno nuevo: cómo gravar el volumen creciente de transacciones entre empresas dentro de cadenas de producción cada vez más largas y complejas. La respuesta no fue homogenea en todas partes. 

Alemania y Francia adoptaron, de forma paralela, un impuesto general sobre el volumen de negocios, que gravaba cada venta sucesiva sin ningún mecanismo de deducción del impuesto ya soportado en la compra anterior. Diseño que resucitaba exactamente el mismo problema de acumulación piramidal que ya habían sufrido las aduanas internas francesas del Antiguo Régimen, solo que ahora entre fases sucesivas de una misma cadena productiva en lugar de entre provincias. 

\paragraph{Impuesto sobre el valor Añadido: LAURÉ (1953)}
La solución llegó de la mano de Maurice LAURÉ, subdirector de la recién creada Dirección General de Impuestos francesa, quien formuló en 1953 y vio aprobada por la Asamblea Nacional el 10 de abril de 1954 la \textit{taxe sur la valeur ajoutée}, tras haberla ensayado primero, como experimento piloto, en la colonia francesa de Costa de Marfil.\footnote{%
    El contexto político en que esta reforma se aprueba tampoco es casual. Francia atravesaba entonces una intensa ola de protesta fiscal popular, liderada por Pierre Poujade y su Unión de Defensa de Comerciantes y Artesanos, que denunciaba el sistema de imposición en cascada como fuente de arbitrariedad e injusticia entre pequeños y grandes negocios.
}

Lo que resulta interesante, desde nuestra óptca, son los tres objetivos de la reforma que el propio LAURÉ explicitó, ya que son los mismos tres criterios que cualquier propuesta de reforma sigue empleando para evaluar un impuesto sobre transacciones: (i) asegurar un nivel de \textbf{recaudación estable}; (ii) mantener la \textbf{neutralidad} frente a las decisiones de inversión y de exportación de las empresas, el problema exacto que el impuesto en cascada distorsionaba; y, (iii) \textbf{minimizar la vulnerabilidad al fraude}, precisamente mediante el mecanismo de crédito de factura que obliga a cada empresa de la cadena a documentar tanto sus ventas como sus compras.

\paragraph{Impuesto sobre el Valor Añadido: Proceso de expansión}
El Tratado de Roma de 1957 sentó las condiciones para que, el 11 de abril de 1967, el Consejo aprobara la Primera y la Segunda Directivas IVA, obligando a todos los Estados miembros a sustituir sus impuestos nacionales sobre el volumen de negocios (o análogos) por un sistema común basado en el modelo de LAURÉ.

Desde entonces, el modelo se ha extendido a más de $170$ países de los cinco continentes, convirtiéndose, en apenas siete décadas, en la figura de imposición general sobre el consumo dominante a escala prácticamente universal; ritmo que no tiene parangón en ninguna otra gran figura tributaria de la historia contemporánea. De hecho, la excepción más señalada a esta convergencia global es Estados Unidos: la única gran economía del mundo que, pese a haber discutido en varias ocasiones propuestas serias de adopción, nunca ha llegado a introducir un IVA de alcance federal, manteniendo en su lugar el sistema de impuestos estatales sobre ventas minoristas de diseño monofásico que hereda, en última instancia, la misma lógica del punto único de control que ya vimos operar en las aduanas romanas y en las gabelas del Antiguo Régimen.\footnote{%
    Un cierre de círculo que conviene tener presente: la figura más antigua y la ausencia más señalada comparten, en el fondo, el mismo principio elemental de diseño.
}

\subsection{Argumentos a favor: ¿cómo justificamos la imposición indirecta?}
Ahora que conocemos, en términos generales, los grandes hitos en la historia de las figuras de imposición indirecta, cabe preguntarse la otra pregunta crucial: ¿por qué el Estado interviene sobre el acto de gastar la riqueza, en lugar de limitarse a intervenir sobre su generación o su acumulación? No es una pregunta retórica. Gravar una transacción exige un fundamento propio, distinto del que legitima gravar una renta o un patrimonio, porque el hecho que se grava es de otra naturaleza: no es un estado de cosas que persiste en el tiempo, sino un acto puntual e irrepetible, el momento exacto en que un bien o un servicio cambia de titularidad a cambio de una contraprestación. 

Dicho de otra forma, examinar la historia nos permite entender por qué las figuras de tributación indirecta tienen un recorrido histórico mucho mayor al de las figuras directas. Pero la justificación normativa es una cuestión que, directa o indirectamente, tenía nula cabida durante toda ella: ¿qué autócrata siente la necesidad de justificar sus decisiones? ¿qué familia medieval debía justificar la tributación en términos de equidad o eficiencia? Estas preguntas, absurdas en el contexto que hemos presentado, cobran un sentido especial en el Estado moderno: no se puede actuar sin más, el gobernante debe justificar sus decisiones y, en especial, aquellas que inciden sobre derechos que se consideran inalienables/humanos. 

Existen varias justificaciones al respecto. 

\subsubsection{Equidad: Imposición sobre el gasto (HOBBES, 1651; RAWLS, 1971)}
La más antigua y filosóficamente más elaborada a favor de gravar el consumo en lugar de la riqueza no procede de ningún economista, sino de HOBBES (Leviatán, 1651), en su famosa obra, escrita en plena Guerra Civil inglesa, un cuarto de siglo antes de que naciera siquiera la disciplina que hoy llamamos economía política. 

Su formulación exacta merece citarse porque cada palabra está calculada: \textit{la igualdad de la imposición consiste más en la igualdad de lo que se consume, que en la de las riquezas de las personas que lo consumen}. El razonamiento que sostiene esta frase parte de una pregunta previa: ¿qué debe cada súbdito al Estado a cambio de la protección que este le proporciona? Para HOBBES, ese deber no se mide por lo que una persona posee, sino por lo que efectivamente extrae del conjunto de bienes y servicios que la sociedad, bajo la protección del soberano, pone a su disposición. 

Imagina, siguiendo la lógica del propio HOBBES, a dos hombres con idéntica riqueza acumulada: uno vive con frugalidad, ahorra la mayor parte de sus ingresos y apenas consume; el otro gasta sin freno, disfrutando de todo lo que su riqueza le permite. Bajo un impuesto sobre la riqueza o la renta, ambos pagarían lo mismo, porque ambos tienen la misma capacidad económica real. Pero, el segundo ha \textit{tomado} mucho más del acervo común de bienes disponibles en la sociedad que el primero: ha ocupado más recursos, ha empleado más trabajo ajeno, ha disfrutado de más protección efectiva del Estado sobre los bienes que ha llegado a consumir; mientras que el primero, por el simple hecho de no haber gastado, no le debe a la comunidad nada más allá de lo que ya le ha devuelto en forma de ahorro disponible para generaciones futuras o para la propia economía. La equidad, para HOBBES, no reside en igualar cargas sobre la posesión, sino en igualar cargas sobre la extracción efectiva.

Este argumento, ignorado durante siglos por una doctrina hacendística que se construyó casi enteramente sobre la capacidad de pago, encontró una defensa inesperada tres siglos después. En \textit{A Theory of Justice} (1971), el filósofo Jhon RAWLS sostiene que, bajo instituciones de fondo justas y con una renta obtenida de forma legítima, un impuesto proporcional sobre el gasto sería preferible a un impuesto progresivo sobre la renta, y ofrece dos razones que son, en esencia, la reformulación exacta del argumento de HOBBES En primer lugar, que un impuesto sobre el gasto grava en función de cuánto toma una persona del acervo común de bienes, y no en función de cuánto contribuye a él. La segunda, es que trata a todos de manera uniforme, en el sentido de que no penaliza fiscalmente la decisión de ahorrar y diferir el consumo. 

En esencia, los argumentos de HOBBES y RAWLS muestran que la elección entre gravar la renta o el gasto no se alinea de forma automática con ningún eje ideológico convencional, sino que responde a una pregunta filosófica autónoma sobre qué mide realmente la justicia distributiva: la posesión, la contribución, o la extracción. Pregunta que, como sabemos, aún no se ha resuelto definitivamente. 

\subsubsection{Equidad: Renta informal}
Junto a este argumento de equidad filosófica, la justificación contemporánea añade uno de naturaleza puramente instrumental, pero con un respaldo empírico que conviene traer con datos concretos y no solo como intuición razonable. 

La imposición sobre el consumo funciona como una segunda oportunidad de observación para el Estado. Toda renta que haya escapado a la verificación en el momento de generarse (trabajo informal, evasión de IVA, etc.) termina, tarde o temprano, materializándose en un acto de gasto, momento en el que vuelve a ser, al menos potencialmente, observable.

La magnitud de lo que hay en juego con este argumento no es en absoluto trivial. SCHNEIDER (2007), una de las estimaciones más completas sobre el tamaño de la economía sumergida a escala mundial, sitúan el tamaño medio de la economía no declarada en torno al $16,3\%$ del PIB en los países de la OCDE, y hasta el $39,1\%$ en los países en desarrollo y el $40,1\%$ en las economías en transición. Como vemos, aunque sea una estimación y esté sujeta a crítica, estas cifras aproximan la medida real de cuánta actividad económica escapa, de entrada, a cualquier forma de imposición directa sobre la renta, y explican por qué ningún sistema fiscal comparado ha considerado nunca razonable prescindir de una base tributaria alternativa capaz de captar, aunque sea parcialmente, esa actividad. Ahora bien, es evidente que esta justificación teórica no es perfecta. Evidentemente solo funciona cuando el consumo financiado con renta no declarada se realiza, a su vez, dentro del circuito formal de la economía. Si la propia transacción de consumo se produce también dentro de la economía sumergida, la cadena de observación se rompe en ambos extremos a la vez, y ni la renta ni el gasto llegan nunca a ser visibles para la Administración. 

La imposición sobre el consumo reduce, por tanto, el margen de renta que puede permanecer indefinidamente invisible, pero no lo elimina.

\subsection{Argumentos en contra: ¿por qué no deberíamos tener impuestos indirectos?}
Cierro el Bloque 0 con el contrapunto que, siguiendo el mismo patrón que ya aplicamos a la imposición directa, no puede faltar: los argumentos que, desde la propia tradición clásica de la hacienda pública, han cuestionado —no la conveniencia recaudatoria de gravar el consumo, que nadie discute en serio, sino su legitimidad democrática, la forma en que ese gravamen se articula. No traigo aquí la crítica de la regresividad, que desarrollaremos con todo su aparato técnico más adelante: las que siguen son objeciones de naturaleza distinta, dirigidas no a quién soporta la carga sino a si el propio contribuyente sabe que la está soportando, y qué consecuencias tiene para el control democrático del gasto público que no lo sepa.

\subsubsection{Eficiencia: Principio de la transparencia en la recaudación-gasto (MILL, 1848)}
Resulta paradójico que la crítica más influyente a la imposición indirecta proceda del mismo autor que formalizó, en 1848, la propia distinción entre imposición directa e indirecta que veremos más adelante. 

MILL (1848) sostiene que un ciudadano debería poder conocer, con la máxima claridad posible, cuánto contribuye exactamente al sostenimiento del Estado, porque solo esa conciencia clara de la propia carga fiscal le permite ejercer un control efectivo sobre cómo se gasta lo recaudado. En este sentido, los impuestos indirectos, al quedar incorporados en el precio final del bien o servicio, violan sistemáticamente esta condición: el consumidor paga el impuesto en cada compra sin que, salvo que examine deliberadamente el desglose de la factura, tenga conciencia clara de cuánto de ese precio corresponde al bien y cuánto al tributo. MILL formula, en consecuencia, una preferencia normativa fuerte por la imposición directa, no por razones de eficiencia económica ni de distribución de la carga, sino por una razón estrictamente democrática: un sistema fiscal opaco erosiona la capacidad del ciudadano para actuar como contrapeso efectivo frente al poder de gasto del Estado.\footnote{%
    El economista italiano PUVIANI (1903) dió un paso mucho más radical sobre este mismo terreno intelectual. No se trata de que la imposición indirecta tienda a ocultar la carga fiscal como efecto colateral no buscado, sino de que el Estado elige deliberadamente los instrumentos fiscales que generan mayor ilusión, precisamente porque una menor percepción del coste real del gobierno facilita la aceptación social de un nivel de gasto público más elevado del que los ciudadanos tolerarían si conocieran su coste exacto. PUVIANI clasifica sistemáticamente, en un catálogo detallado que ocupa buena parte de su obra, las distintas técnicas mediante las cuales el poder público oculta tanto el coste de los tributos (ilusiones en los ingresos) como sobreestima el beneficio del gasto que financian (ilusiones en el gasto). Evidentemente, la imposición indirecta, incorporada en el precio y fraccionada en miles de pequeños pagos cotidianos, ocupa el lugar central de su catálogo: el instrumento más eficaz jamás diseñado para este fin.

La obra de PUVIANI permaneció prácticamente ignorada durante más de medio siglo, hasta que BUCHANAN la redescubrió en los años sesenta y la incorporó como uno de los pilares de su propio programa de investigación sobre las restricciones constitucionales al poder fiscal del Estado, el mismo terreno del que surgiría después la crítica del Estado Leviatán junto a BRENNAN. La conexión no es casual: si el Leviatán de BRENNAN y BUCHANAN es un Estado que explota al máximo cualquier base imponible amplia que se le conceda, la teoría de la ilusión financiera de PUVIANI explica el mecanismo concreto mediante el cual ese Leviatán consigue, además, que esa explotación genere la mínima resistencia política posible.
}

Este argumento, lejos de quedar como una curiosidad decimonónica, ha generado un programa de investigación empírica propio y activo hasta la actualidad. SAUSGRUBER y TYRAN (2005), entre otros, diseñaron experimentos de laboratorio para contrastar directamente si los sujetos perciben de forma sistemáticamente distinta cargas fiscales idénticas en su magnitud pero distintas en su grado de visibilidad. Como cabía esperar, encontraron evidencia consistente con la propia intuición de MILL: la forma de presentar el impuesto altera, de manera significativa y medible, cómo lo perciben quienes lo pagan, con independencia de que la carga económica real sea idéntica en ambos casos.\footnote{%
    Por su parte, CHETTY, LOONEY y KROFT (2009) realizaron una de las pruebas directas más sobresalientes al respecto. Los autores diseñaron un experimento de campo en un supermercado real, en el que sustituyeron las etiquetas de $750$ productos sujetos a impuesto sobre ventas por etiquetas que mostraban explícitamente el precio final incluido el impuesto, en lugar del precio sin impuesto habitual (en Estados Unidos, el impuesto se añade en caja). El resultado fue una caída de la demanda de los productos tratados de aproximadamente un $8\%$, sin que hubiera cambiado en absoluto el precio real pagado en caja: el único cambio fue la percepción del consumidor. Un refinamiento posterior del propio experimento mostró, además, que una parte no desdeñable de esa caída de demanda no se debía a la mera falta de atención al precio final, sino a un problema de información aún más básico: cerca de un $20\%$ de los compradores encuestados ni siquiera sabía que un producto tan cotidiano como la pasta de dientes estaba sujeto a impuesto alguno.
}

\subsubsection{Conductual: Incentivos perversos (SMITH, 1776)}
Cierro con una tercera crítica, de naturaleza distinta a las dos anteriores porque no versa sobre la percepción del contribuyente sino sobre el coste efectivo de administración, y que procede de SMITH (1776). En el mismo pasaje donde formulaba sus cuatro cánones de la imposición, dedica el cuarto de ellos a advertir que un buen impuesto debe recaudarse sacando y reteniendo de los bolsillos del pueblo lo menos posible, por encima de lo que aporta al tesoro público. 

SMITH aplicó este criterio con particular dureza a los aranceles y a los impuestos especiales de su época, precisamente los antecesores directos de la imposición indirecta moderna. Su objeción no era de principio sino práctica, y con tres vertientes que conviene distinguir. Primera, el propio coste del aparato de vigilancia necesario para impedir el contrabando y el fraude podía llegar a consumir una parte muy sustancial de lo efectivamente recaudado. Segunda, el propio incentivo estructural al contrabando que genera cualquier tipo elevado sobre un bien fácilmente transportable, con el coste social añadido de criminalizar a comerciantes que, en ausencia del impuesto, no habrían incurrido en ninguna conducta ilícita. Es decir, para SMITH, un impuesto indirecto excesivamente alto podía convertir a \textit{un hombre que, en todos los demás aspectos, sería un buen ciudadano} en un contrabandista, simplemente por la magnitud del incentivo económico creado por el propio Estado. Tercera, la distorsión que la persecución del fraude introduce sobre la propia actividad comercial legítima, al someterla a inspecciones, registros y trabas administrativas que penalizan también a quien no tiene ninguna intención de defraudar.

En definitiva, MILL advierte del coste que su opacidad impone sobre la responsabilidad democrática; PUVIANI sugiere que esa opacidad no es un defecto accidental sino, potencialmente, una característica buscada por el propio poder que la administra; y SMITH recuerda que, incluso dejando de lado cualquier consideración de transparencia, el propio acto de hacer cumplir un impuesto sobre transacciones tiene un coste de ejecución que ningún diseño técnico, por sofisticado que sea, puede reducir por completo a cero. 





\clearpage
\section{Objeto: ¿?}
\puntosclave{
    
}

Fijado ya el fundamento, corresponde ahora precisar qué se grava exactamente cuando se grava el consumo. No es una pregunta tan sencilla como parece a primera vista: un mismo bien físico puede pasar por múltiples manos antes de llegar al consumidor final y el diseño técnico elegido para gravar esa cadena tiene consecuencias que van mucho más allá de la mera recaudación: determina si el sistema penaliza o no la organización empresarial en varias fases sucesivas, si genera incentivos a la integración vertical artificial, y si resulta o no vulnerable a determinadas formas de fraude. 

Vamos ahora a recorrer cuatro preguntas sobre el objeto de la imposición indirecta: (i) qué diseño estructural adoptar frente a una cadena de producción con múltiples fases; (ii) cómo calcular, técnicamente, la magnitud exacta que se pretende gravar en cada fase; (iii) si ese gravamen debe aplicarse en el lugar donde el bien se produce o en el lugar donde se consume, cuando ambos no coinciden; y, (iv) si debe aplicarse de forma general sobre cualquier transacción o de forma selectiva sobre determinados bienes. 

\subsection{Diseño: ¿cómo diseñar la arquitectura del impuesto?}


\subsubsection{Gravar la venta: ¿por qué la fase importa?}

\paragraph{Monofásica: gravar un único eslabón de la cadena}
La solución más simple, en apariencia, consiste en gravar la transacción en un único punto de la cadena de producción y distribución, dejando el resto libre de gravamen. Pero esta aparente simplicidad esconde una decisión nada trivial: qué eslabón elegir, porque cada opción reparte de forma distinta el equilibrio entre amplitud de la base y número de contribuyentes a administrar. 

Gravar en la fase de fabricación tiene la ventaja de concentrar la obligación tributaria en un número relativamente reducido de grandes productores, pero dejar sin gravar los márgenes de distribución mayorista y minorista que se añaden después; por tanto, implica una base sustancialmente más estrecha que el valor final que paga el consumidor. Por el contrario, gravar en la fase minorista (sales tax estadounidense) captura la totalidad del valor final del bien, porque se aplica justo antes de que llegue al consumidor, pero multiplica exponencialmente el número de sujetos obligados a liquidar el impuesto: mientras que un país puede tener unos pocos miles de fabricantes de un determinado bien, puede tener cientos de miles o millones de puntos de venta minorista, lo que encarece sustancialmente el coste de control y de cumplimiento.

El diseño monofásico retail, además, tiene una fragilidad estructural que la economía digital ha puesto de manifiesto. Si el impuesto se cobra en un único punto, todo el sistema depende de que ese punto sea identificable y esté sujeto a la jurisdicción fiscal correspondiente. El Tribunal Supremo de Estados Unidos (South Dakota v. Wayfair, Inc., 2018), tuvo que resolver exactamente este problema: durante 26 años, ningún Estado podía exigir a un vendedor sin presencia física en su territorio que recaudara su impuesto sobre ventas. Esta doctrina, absolutamente desfasada temporalmente, permitía que, con la explosión del comercio electrónico, gigantes de la comercio electrónico retail vendieran masivamente a consumidores de un Estado sin recaudar nunca el impuesto correspondiente, erosionando la base de cualquier Estado con fuerte penetración de comercio online.\footnote{%
    El Tribunal, por una ajustada mayoría de cinco votos frente a cuatro, revocó expresamente Quill, considerando que generaba distorsiones al comercio interestatal pues protegía artificialmente al vendedor sin presencia física frente al competidor local, y validó el criterio perseguido por la ley de Dakota del Sur que exigía la recaudación a cualquier vendedor, sustituyendo el criterio de presencia física por uno de nexo económico (umbral de ventas o de transacciones).
}

Por tanto, vemos como un problema del diseño monofásico es que, cuando el punto único de control deja de ser fácilmente localizable, todo el sistema de recaudación puede colapsar por su propio diseño.

\paragraph{Plurifásico: Acumulación piramidal}
La alternativa histórica al diseño monofásico es gravar todas las fases de la cadena: el impuesto plurifásico. Este impuesto con contempla ningún mecanismo de deducción del impuesto ya soportado en la compra anterior, por tanto constituye un impuesto en cascada o sobre el volumen de negocios. 

Conviene detenerse en la naturaleza exacta del problema que esto genera, porque no es simplemente se pague más impuesto. Si cada una de las $n$ fases de una cadena de producción aplica un tipo $t$ sobre el precio ya gravado en la fase anterior, la carga acumulada final no crece de forma lineal con el número de fases, sino de forma compuesta, exactamente igual que un interés que se capitaliza sobre sí mismo: el tipo efectivo final sobre el valor añadido total se aproxima a $(1+t)^n− 1$, no a $n\cdot t$. Esto significa que una cadena de producción con cinco fases sucesivas, cada una gravada al $10\%$, no soporta una carga acumulada del $50\%$, sino sensiblemente superior, precisamente porque cada fase paga impuesto no solo sobre su propio valor añadido, sino sobre todo el impuesto ya pagado en las fases anteriores. 

La consecuencia económica de este diseño es una distorsión estructural sobre la propia organización empresarial: una empresa que consiga integrarse verticalmente, eliminando transacciones intermedias sujetas a gravamen, reduce su carga fiscal total frente a un competidor que mantenga una cadena fragmentada en varias empresas independientes, con exactamente el mismo valor añadido total. Este incentivo, documentado en la historia económica alemana durante el periodo de entreguerras, contribuyó de forma nada desdeñable a la concentración empresarial en grandes conglomerados verticalmente integrados en la Alemania de los años veinte y treinta. Por tanto, un efecto colateral del diseño fiscal casi inevitable es sobre la propia estructura de mercado.

\subsubsection{Gravar el valor añadido: ¿por qué la fase deja de importar?}
El tercer diseño resuelve el problema de raíz mediante un mecanismo que ilustra mejor un ejemplo. Imaginemos una cadena de tres fases. Un fabricante que vende a un mayorista por $100$, el mayorista que vende al minorista por $150$, y el minorista que vende al consumidor final por $200$, con un tipo del $20\%$ en todas las fases. 

Bajo el crédito de factura, diseñado por LAURÉ, el fabricante repercute $20€$ y, como no ha soportado impuesto alguno en sus propias compras, ingresa esos $20€$ íntegros a la Hacienda Pública. El mayorista repercute $30€$ en su venta, pero deduce los $20€$ que soportó en su compra al fabricante, ingresando solo la diferencia, $10€$, el $20\%$ de su propio valor añadido. El minorista repercute $40€$ y deduce los $30€$ que soportó, ingresando $10$, el $20\%$ de su propio valor añadido. La suma de lo ingresado por las tres empresas a lo largo de toda la cadena es $20 + 10 + 10 = 40$, exactamente el $20\%$ del precio final de $200€$ que paga el consumidor, con independencia total de en cuántas fases se haya fragmentado la cadena. 

Por tanto, por su formulación, mientras que en la cascada el número de fases determinaba directamente la carga total, en el crédito de factura el número de fases es, para la recaudación final, completamente irrelevante: solo importa el valor añadido total acumulado hasta el consumidor, exactamente la magnitud que el impuesto pretende gravar.

\paragraph{Teorema de eficiencia productiva: DIAMOND y MIRRLEES (1971)}
LAURÉ llegó a este diseño por pura necesidad práctica, para resolver el problema de la acumulación piramidal que ya atormentaba a la industria francesa de posguerra. Pero existe un argumento teórico que demuestra porque, además de resolver un tormento administrativo e industrial, esta arquitectura es óptima en términos económicos.

No obstante, esta formulación llegaría casi veinte años después, de la mano de DIAMOND y MIRRLEES (1971). Los autores demuestran, dentro del marco general de la teoría de la imposición óptima (\authorfont{Ver Tema 3.B.11}), el llamado Teorema de eficiencia productiva: \textbf{bajo cualquier sistema óptimo de imposición, ningún bien intermedio debe tributar, con independencia de que el resto del sistema fiscal esté plagado de impuestos distorsionantes sobre otras magnitudes}.\footnote{%
    \textcolor{red}{\textbf{NOTA.-} Es muy importante incluir una breve demostración (gráfica o analítica del Teorema)}
} 

La intuición es exactamente la que el ejemplo numérico anterior acaba de mostrar de forma mecánica. Gravar una transacción entre empresas distorsiona las decisiones de producción, ya que incide sobre qué insumos emplear, cómo organizar la cadena, si integrar verticalmente o subcontratar. Además, esta distorsión no genera a cambio ningún beneficio de bienestar social, porque el bienestar depende únicamente del consumo final, nunca de cómo se organiza internamente la producción para llegar a él. Por tanto, toda la carga tributaria debe concentrarse en el punto en que el bien llega efectivamente al consumidor, dejando intacta la asignación de factores entre las empresas intermedias (eficiente). 

El crédito de factura, con su propiedad de neutralidad ante el número de fases es exactamente la traducción práctica de esa condición de eficiencia: cada fase intermedia soporta, en términos netos, una carga de cero sobre las transacciones que mantiene con otras empresas, y toda la recaudación efectiva recae, sin excepción, sobre el margen final que separa al último eslabón de la cadena del consumidor. 

\subsection{Cálculo de la base: ¿cuál es la forma óptima?}


\subsubsection{Métodos de cálculo: ¿cómo calcular el valor añadido?}
Antes de asumir, como hace la práctica casi universal, que el valor añadido de una empresa solo puede calcularse mediante el crédito de factura que acabamos de desarrollar, conviene detenerse en algo que rara vez se explicita con la precisión que merece: existen, en teoría, tres métodos matemáticamente equivalentes de llegar exactamente al mismo resultado, y el hecho de que el mundo haya convergido casi por completo en uno solo de ellos no es una necesidad lógica, sino una decisión de diseño con razones concretas que conviene entender una por una.

\paragraph{Sustracción: Agregación de ventas y compras}
El primero es el método de sustracción. El valor añadido de una empresa en un periodo se calcula, directamente, como la diferencia entre sus ventas totales y sus compras totales a otras empresas, sin necesidad de descomponer esa cifra transacción por transacción. Por tanto, el impuesto a ingresar es, simplemente, el tipo aplicado sobre esa diferencia agregada, calculada al cierre del periodo a partir de la contabilidad general de la empresa, de forma análoga a como se calcula el beneficio contable.

\paragraph{Adición: Agregación del total remunerado a factores}
El segundo es el método de adición. Dado que, por la propia definición contable de valor añadido, este debe coincidir exactamente con la suma de las retribuciones pagadas a los factores de producción empleados en esa fase, el valor añadido puede calcularse también sumando estas partidas directamente desde la cuenta de pérdidas y ganancias de la empresa, sin necesidad de referencia alguna a las cifras de venta y compra.

\paragraph{Crédito de factura: Desagregar por transacción}
El tercero es el crédito de factura. En este caso, no se calcula el valor añadido como magnitud agregada en ningún momento, sino que se calcula directamente el impuesto a ingresar como la diferencia entre el impuesto repercutido en las ventas y el impuesto soportado en las compras, transacción por transacción.

\subsubsection{Teorema de equivalencia: ¿por qué todo el mundo converge en el tercer método?}
Los tres métodos son, bajo el supuesto de que se aplique un tipo único y uniforme a toda la economía, sin exenciones ni tipos reducidos, exactamente equivalentes en su resultado matemático.\footnote{%
    \textcolor{red}{\textbf{NOTA.-} Es importante incluir una demostración breve de esto.}
} Cualquiera de los tres, aplicado correctamente, arroja idéntica recaudación total y distribuye idéntica carga entre las fases de la cadena. Es, en el fondo, la misma propiedad contable elemental que asegura que el beneficio de una empresa puede calcularse, indistintamente, restando los gastos de los ingresos o sumando las partidas que componen ese beneficio desde otra perspectiva.

Si los tres métodos son equivalentes, ¿por qué el mundo ha convergido de forma tan abrumadora en el crédito de factura, relegando a la sustracción y la adición a un papel casi anecdótico? 

\paragraph{Problemas de implementación: Diferenciación de tipos}
La respuesta no está en ninguna superioridad teórica del tercer método, sino en que ese supuesto casi nunca se cumple el supuesto que exige que esta equivalencia se materialice. Como veremos, la inmensa mayoría de los sistemas de imposición sobre el consumo aplican tipos diferenciados según el tipo de bien y, por tanto, los métodos de sustracción y de adición no arrojan el mismo resultado, ya que calculan el valor añadido agregado al cierre del periodo a partir de la contabilidad general de la empresa y sin ninguna referencia a qué tipo concreto de bien generó cada venta o cada compra individual. Si una empresa vende, en el mismo periodo, productos sujetos al tipo general y productos sujetos a un tipo reducido, el método de sustracción, aplicado sobre la cifra agregada de ventas menos compras, no tiene forma de saber qué proporción de esa diferencia corresponde a qué tipo.\footnote{%
    Salvo que se lleve, de facto, una contabilidad separada por tipo de bien. Lo que en realidad equivale, en la práctica, a reconstruir la información que el crédito de factura ya proporciona de forma nativa, transacción por transacción
}

Existe un experimento natural extraordinariamente reciente que confirma ésto con. Japón introdujo su impuesto sobre el consumo en abril de 1989, y durante treinta años, hasta 2019, lo administró mediante un sistema mucho más próximo al método de sustracción que al crédito de factura europeo. Los sujetos pasivos podían deducir el impuesto soportado en sus compras basándose simplemente en sus libros y registros contables propios, sin necesidad de conservar ni verificar facturas formales expedidas por sus proveedores. Este diseño funcionaba razonablemente bien precisamente porque, durante esas tres décadas, Japón mantuvo un tipo único para la totalidad de bienes y servicios, sin ninguna diferenciación. 

Sin embargo, el problema se hizo evidente cuando Japón rompió esa uniformidad. En octubre de 2019, al elevar el tipo general del $8\%$ al $10\%$, mantuvo un tipo reducido del $8\%$ para alimentos y bebidas. El sistema basado en libros contables, sin factura formal, resultó inmediatamente insuficiente. Numerosos sujetos pasivos empezaron a computar de forma incorrecta su crédito fiscal, precisamente por no poder acreditar con precisión, a partir de su propia contabilidad agregada, qué parte de sus compras había soportado el $8\%$ y cuál el $10\%$. Japón introdujo primero, como medida transitoria desde 2019, un sistema de facturas clasificadas por tipo que exigía que las facturas indicaran expresamente el tipo aplicable, y finalmente, desde 2023, implantó el crédito de factura como sistema uniforma, que exige facturas formales con número de registro del emisor, desglose del importe y de la cuota por cada tipo aplicable, y establece que las entidades no registradas como sujetos pasivos cualificados no pueden expedir facturas con derecho a deducción, de modo que sus clientes pierden el derecho a deducir el impuesto correspondiente a esas compras.

En apenas cuatro años, Japón se vio obligado a abandonar un sistema de sustracción que había funcionado sin sobresaltos durante tres décadas, y a converger hacia un diseño prácticamente indistinguible del crédito de factura europeo. Fenómeno que prueba que la elección del método de cálculo no es una cuestión de preferencia, sino una consecuencia necesaria de la decisión tomada en otro terreno completamente distinto: renunciar a la uniformidad tributaria.

\paragraph{Incentivos: Autocontrol documental}
Existe una segunda ventaja del crédito de factura, distinta de la compatibilidad con tipos diferenciados, que conviene no perder de vista porque conecta directamente con el problema del fraude \textcolor{red}{que ya analizamos en el Bloque 0}. 

Al exigir que cada empresa documente, mediante factura, tanto sus ventas como sus compras para poder ejercer su derecho a deducción, el sistema genera un incentivo espontáneo a que cada comprador exija a su proveedor una factura válida, lo que crea una cadena de verificación cruzada entre empresas independientes, cada una vigilando, por puro interés propio, que la anterior haya documentado correctamente la operación. Ni el método de sustracción ni el de adición generan este efecto, porque ambos se basan en la contabilidad interna y agregada de cada empresa, sin ninguna referencia obligada a lo que documenta la contraparte de cada transacción. Es, precisamente, esta propiedad de autocontrol la que LAURÉ identificaba como uno de sus tres objetivos de diseño: minimizar la vulnerabilidad al fraude, y la que, paradójicamente, sigue siendo hoy el terreno de la reforma más activa del sistema en toda Europa. Por ejemplo, la iniciativa \textit{VAT in the Digital Age} (ViDA, 2024) de la Comisión, generalizará la facturación electrónica obligatoria y la comunicación digital en tiempo casi real de cada transacción sujeta a IVA transfronterizo a partir de 2030, llevando la lógica del autocontrol documental a su forma más avanzada: la verificación automatizada de cada factura por parte de la propia Administración en el momento mismo en que se emite.\footnote{%
    Esta iniciativa viene precedida por experiencias nacionales pioneras como el sistema italiano de facturación electrónica obligatoria (2019) o el Suministro Inmediato de Información español.
}

\subsection{Destino frente a origen: ¿qué país tiene derecho a gravar una transacción que cruza una frontera?}
Todo lo desarrollado hasta ahora en este bloque presuponía, de forma implícita, que productor y consumidor de un bien se encuentran dentro de la misma jurisdicción fiscal. En cuanto un bien o un servicio cruza una frontera, surge una pregunta que no se resuelve por sí sola: ¿debe tributar en el país donde se produjo o en el país donde finalmente se consume? 

El principio de origen grava la transacción en el país productor, con el tipo vigente en ese país, y el bien viaja después con ese impuesto ya incorporado, sin ajuste adicional al cruzar la frontera. El principio de destino hace exactamente lo contrario: el país exportador exonera la venta (desgravación en frontera), devolviendo o no cobrando el impuesto que de otro modo correspondería, y el país importador aplica su propio tipo en el momento en que el bien entra en su territorio, de modo que la carga fiscal final depende de dónde se consume.

\subsubsection{Neutralidad competitiva: ¿por qué el mundo ha convergido en el principio de destino?}
La razón por la que el mundo ha convergido, de forma casi unánime, en el principio de destino es una razón de neutralidad competitiva dentro de cualquier mercado nacional dado. Bajo el principio de origen, un bien fabricado en un país con tipo del $25\%$ competiría, en cualquier mercado del mundo, en desventaja permanente frente a un bien idéntico fabricado en un país con tipo del $10\%$, con independencia de la eficiencia productiva real de cada fabricante: la ventaja competitiva dependería del sistema fiscal del país de origen, no de la productividad (como mínimo parcialmente). Bajo el principio de destino, en cambio, dos bienes idénticos que compiten dentro del mismo mercado de consumo soportan exactamente la misma carga fiscal, con independencia de dónde se hayan producido, preservando así la neutralidad competitiva dentro de cada mercado nacional.

\paragraph{Organización Mundial del Comercio: Ajuste en frontera del IVA}
Esta preferencia teórica por el principio de destino tiene, además, un respaldo jurídico internacional que conviene conocer porque revela una asimetría poco conocida y con enormes consecuencias prácticas.

Las normas de la OMC permiten el ajuste fiscal en frontera para los impuestos indirectos como el IVA, desgravar las exportaciones y gravar las importaciones al tipo doméstico se considera una simple corrección técnica, no una ventaja comercial desleal. Sin embargo, las mismas normas prohíben ese mismo ajuste para los impuestos directos. Es decir, un país no puede devolver a sus exportadores el impuesto sobre sociedades que han pagado sobre el beneficio generado por sus ventas al exterior, porque eso se consideraría una subvención a la exportación ilegal bajo las normas comerciales internacionales.\footnote{%
    Esta distinción, aparentemente jurídica, fue el epicentro de uno de los debates más intensos de la última década: la propuesta de impuesto sobre el flujo de caja ajustado en destino que se defendía como sustituto del impuesto sobre sociedades tradicional, proponiendo, precisamente, ajustar en frontera un impuesto directo, exonerando las exportaciones y gravando las importaciones. La propuesta se topó con una oposición furibunda de grandes minoristas fuertemente dependientes de importaciones, además de con serias dudas sobre su compatibilidad con las normas de la OMC precisamente por tratarse de un impuesto directo y no indirecto. Finalmente, terminó siendo abandonada, aunque el episodio ilustra que la frontera entre lo que el DIPr permite ajustar en frontera y lo que prohíbe no es un detalle técnico menor.
}

\paragraph{Parche temporal: ¿qué criterio aplica la Unión en las transacciones del MIU?}
El episodio que mejor ilustra la dificultad práctica de esta disyuntiva es el de la propia Unión con el mercado interior. 

Cuando la Comunidad Europea preparaba la creación del mercado único, en 1993, el plan original era ambicioso: establecer, para esa misma fecha, un sistema definitivo basado en el principio de origen para las transacciones intracomunitarias, coherente con la lógica de que un verdadero mercado único, sin fronteras interiores, no debería necesitar ajustes fiscales en unas fronteras que, formalmente, ya no existían. Sin embargo, el plan resultó inviable por la falta de armonización de los tipos de IVA entre Estados miembros y, en consecuencia, la necesidad de establecer un mecanismo de redistribución de la recaudación entre el país productor, que cobraría el impuesto, y el país consumidor, que se quedaría sin esa recaudación pese a ser donde el bien se consume realmente. La necesidad de establecer este mecanismo, hizo políticamente imposible cerrar el acuerdo a tiempo y se adoptó una solución transitoria mientras se resolvían esas dificultades. ¿Cuál? Precisamente mantener de facto el principio de destino para las transacciones entre empresas: la venta se declara exenta en el país de origen y el comprador declara e ingresa el impuesto correspondiente en el país de destino mediante inversión del sujeto pasivo. Ese régimen transitorio sigue vigente hoy, más de tres décadas después de su implantación.

\textcolor{red}{—y es, precisamente, el mecanismo que ya identificamos en el Bloque 0 como la puerta de entrada exacta del fraude del operador desaparecido: la exención en origen permite a una red criminal adquirir mercancía libre de IVA, revenderla con IVA repercutido en el país de destino, y desaparecer sin ingresar ese impuesto—.} 

La Comisión Europea, consciente de esta relación causal directa entre el parche temporal y las decenas de miles de millones de euros de fraude anual que ya cuantificamos, propuso un sistema definitivo (2019) que invierte por completo el plan original de 1985, abandonando cualquier aspiración al principio de origen y proponiendo, en su lugar, reforzar todavía más el principio de destino; obligando al propio proveedor a repercutir directamente el IVA del país de destino, eliminando la exención que hoy explota el fraude.

Sin embargo, incluso esta propuesta correctiva ha corrido una suerte muy similar. El 11 de febrero de 2025, la Comisión confirmó la retirada definitiva de esta propuesta de sistema definitivo, ante la resistencia de varios Estados miembros preocupados por los nuevos riesgos de fraude y los problemas de recaudación que el nuevo diseño podría generar, desplazando todo el esfuerzo reformista hacia la digitalización de la información fiscal, iniciativa que ya presentamos. Más de tres décadas después de su implantación como solución transitoria, el régimen sigue plenamente vigente sin que ninguna de las dos alternativas teóricas (origen ni el destino reforzado) haya conseguido finalmente sustituirlo. 

\subsection{Diseño del sistema: ¿cuál debería ser la base general y cómo su gravamen?}
La última gran decisión de diseño sobre el objeto de la imposición indirecta será: ¿debe el impuesto aplicarse de forma general, sobre la práctica totalidad de bienes y servicios al mismo tipo, o debe reservarse un tratamiento selectivo para determinados bienes concretos? 

\subsubsection{Amplitud de la base: ¿impuesto general o selectivo?}
La respuesta no es única, porque existen dos fundamentos teóricos completamente distintos que justifican la selectividad, y conviene no confundirlos: uno corrige un fallo de mercado concreto; el otro persigue minimizar la pérdida de eficiencia agregada del sistema fiscal en su conjunto.

\paragraph{Eficiencia: Externalidades negativas (PIGOU, 1920)}
PIGOU (1920) formuló el argumento que sigue siendo el fundamento de referencia de cualquier impuesto selectivo con finalidad correctora: cuando el consumo de un bien genera costes sociales que exceden el coste privado que asume quien lo consume, el mercado produce y consume ese bien en una cantidad superior a la socialmente óptima, porque el precio de mercado refleja solo el coste privado. 

Un impuesto igual, en su cuantía, a ese coste externo marginal, impuesto pigouviano, corrige la distorsión: obliga al consumidor a internalizar en su propia decisión de compra el coste que, de otro modo, trasladaría gratuitamente al resto de la sociedad, restaurando así la eficiencia asignativa sin necesidad de prohibir nada.\footnote{%
    La aplicación contemporánea más estudiada y mejor documentada es el impuesto sobre bebidas azucaradas, introducido el 1 de enero de 2014 —un peso mexicano por litro, en torno al $10\%$ del precio de venta—, en un país que en 2012 tenía más del $70\%$ de su población adulta con sobrepeso u obesidad, y donde las bebidas azucaradas representaban más del 70\% del azúcar añadido en la dieta media. M. Arantxa Colchero y su equipo del Instituto Nacional de Salud Pública mexicano documentó una reducción sostenida de las compras de bebidas gravadas del 5,5\% en el primer año y del 9,7\% en el segundo, acompañada de un incremento del 2,1\% en las compras de bebidas no gravadas, principalmente agua embotellada. El dato más relevante para el propio fundamento pigouviano es que la reducción fue sistemáticamente mayor entre los hogares de menor nivel socioeconómico. Esto es, simultáneamente, la mejor y la peor noticia posible para un impuesto pigouviano: la mejor, porque confirma que el impuesto está modificando efectivamente el comportamiento; la peor, porque revela que la carga económica del impuesto recae también, de forma desproporcionada, sobre los hogares de menor renta.
}

\paragraph{Eficiencia: Minimizar la distorsión (RAMSEY, 1927)}
El segundo fundamento de la selectividad no tiene nada que ver con corregir externalidades: procede de la misma regla de la elasticidad inversa RAMSEY (1927), que encuentra aquí en su formulación original, pensada precisamente para bienes de consumo. 

RAMSEY se plantea un problema puramente de eficiencia: si el Estado necesita recaudar una cantidad fija de ingresos mediante impuestos sobre bienes, y quiere hacerlo minimizando la pérdida de eficiencia agregada, ¿cómo debe repartir la carga entre los distintos bienes disponibles? La pérdida de eficiencia, como sabemos, crece aproximadamente con el cuadrado del tipo aplicado y con la elasticidad de la demanda de ese bien: cuanto más responde la cantidad consumida a una subida de precio, mayor es la cantidad de transacciones mutuamente beneficiosas que el impuesto impide que se realicen. Por tanto, gravar con más intensidad los bienes de demanda inelástica genera mucha menos distorsión que gravar con la misma intensidad bienes de demanda elástica, para una misma recaudación total. 

La regla de RAMSEY concluye, en consecuencia, que el tipo óptimo sobre cada bien debe ser inversamente proporcional a su elasticidad de demanda: tipos más altos sobre lo que menos cambia de comportamiento, tipos más bajos sobre lo que más cambia.\footnote{%
    Conviene resaltar, no obstante, que los bienes de demanda más inelástica tienden a coincidir, en la práctica, con los bienes de primera necesidad. Alimentos básicos, medicamentos, energía doméstica, presentan demandas que varían poco con el precio precisamente porque no son sustituibles ni prescindibles. La eficiencia y equidad, por tanto, tiran en direcciones opuestas. Ningún sistema real sigue la regla de RAMSEY en su forma pura, precisamente por esta tensión.
}

\subsubsection{Formato del impuesto: ¿ad valorem o específico?}
Cerramos con una última pieza técnica, menos conocida fuera del ámbito especializado pero con consecuencias prácticas reales sobre cómo se diseñan estos impuestos selectivos: la elección entre un impuesto ad valorem (un porcentaje sobre el precio de venta) y un impuesto específico (cuantía fija por unidad física). 

\paragraph{Estructura de mercado: Condiciones de equivalencia (SUITS y MUSGRAVE, 1953)}
SUITS y MUSGRAVE (1953) demuestran que, bajo competencia perfecta, ambos formatos son equivalentes para una misma recaudación: da igual cobrar un $10\%$ del precio que una cantidad fija calculada para que, al precio de equilibrio, arroje la misma recaudación. Pero bajo competencia imperfecta esa equivalencia se rompe. En estos mercados, un impuesto ad valorem reduce el margen que el propio productor con poder de mercado obtiene por cada unidad adicional vendida a un precio más alto, lo que debilita su propio incentivo a subir el precio tanto como lo haría en ausencia del impuesto, generando así una subida de precio final y una reducción de cantidad menores que las que provocaría un impuesto específico calculado para recaudar exactamente lo mismo.\footnote{
\textcolor{red}{\textbf{NOTA.-} Conviene aquí demostrar lo expuesto, aunque sea brevemente, pudiendo recurrir tanto al análisis gráfico como al matemático.}
}

La implicación práctica que conviene retener es que la elección entre ad valorem y específico no es un detalle, sino una decisión que interactúa con la propia estructura de mercado del bien gravado. En sectores concentrados y con poder de mercado significativo, el formato ad valorem modera parte del efecto de traslación del impuesto al consumidor, mientras que el formato específico no ofrece esa misma moderación.


\clearpage
\section{Sujeto: }

Ahora que sabemos el objeto sujeto a gravamen, o mejor dicho, la forma en la que podemos diseñar y configurar la base impoible, toca preguntarse otra cuestión: quién liquida el impuesto ante la Hacienda Pública, y si esa persona coincide con aquella a quien la norma pretende, en última instancia, hacer soportar la carga. Vamos a comprobar que, en la imposición indirecta, estas dos figuras casi nunca coinciden por diseño deliberado, y que esa escisión no es un efecto colateral del sistema, sino su característica definitoria más antigua, presente ya en la primera formulación qué distingue a un impuesto indirecto de uno directo. 

Vamos a recorrer cuatro manifestaciones concretas de esa escisión: (i) su formulación conceptual originaria; (ii) su funcionamiento técnico a lo largo de toda una cadena de producción, donde las empresas intermedias actúan, en rigor, como meras recaudadoras; (iii) los umbrales que eximen a ciertos negocios de esa obligación de recaudar por razones de proporcionalidad administrativa; y, (iv) la inversión deliberada de esa misma obligación como respuesta a las formas más sofisticadas de fraude.

\subsection{Sujeto pasivo: ¿tiene sentido hablar de sujeto pasivo en la imposición indirecta?}
El primer autor en proporcionar una distinción comprehensiva y fundada entre la imposición directa e indirecta fue MILL (1848). Sin embargo, el propio vocabulario no nace con él, sino casi un siglo antes, en la Francia de los fisiócratas. 

La escuela de QUESNAY, convencida de que la tierra era la única fuente auténtica de riqueza y de que todo el resto de la actividad económica era estéril, sostenía que cualquier impuesto, con independencia de sobre quién se exigiera formalmente, terminaba recayendo sobre esa única fuente real de riqueza: la renta de la tierra. 

Bajo esa lógica, llamaron directo al único impuesto que gravaba esa renta agrícola sin necesidad de ningún traslado intermedio, e indirecto a cualquier otro, precisamente porque llegaba de forma indirecta hasta el mismo destinatario final. El propio sistema teórico fisiócrata quedó pronto desacreditado y abandonado (impuesto único sobre la tierra), pero, como observó con cierta ironía uno de los economistas de la época, \textit{su sistema ha caído, sus definiciones ya no se admiten, pero sus denominaciones han permanecido en el uso general}. 

El vocabulario, por tanto, sobrevivió al edificio teórico que lo había originado, y es ese vocabulario huérfano de su fundamento original el que MILL recogió, casi un siglo después, para llenarlo de un contenido completamente distinto.

\subsubsection{Distinción: La intención del legislados (Mill, 1848)}
MILL (1848), abandona por completo el criterio fisiócrata y lo sustituye por un criterio radicalmente distinto, basado no en ningún resultado económico observable sino en la intención declarada del legislador. 

Merece la pena citar su formulación exacta, porque cada palabra está calculada con precisión jurídica: un impuesto directo es aquel que se exige de las mismas personas a quienes se pretende o desea que lo paguen. Los impuestos indirectos son aquellos que se exigen de una persona con la expectativa e intención de que esta se resarza a costa de otra: como el impuesto especial o los aranceles. Al productor o importador de una mercancía se le exige pagar el impuesto sobre ella, no con la intención de imponerle a él una contribución particular, sino para gravar, a través de él, a los consumidores de la mercancía, de quienes se supone que recuperará el importe mediante un incremento del precio. Fijémonos en la estructura lógica de esta definición. MILL no está definiendo la naturaleza del bien gravado ni el acto económico que desencadena el impuesto, está definiendo una relación entre dos personas: \textbf{quien la ley obliga a ingresar el tributo, y quien la ley espera e intenta que termine soportándolo mediante un mecanismo de traslación}. En el caso de la imposición indirecta, el incremento del precio. 

Es, en esencia, la primera formulación teórica seria de lo que siglos después la técnica jurídica moderna llamaría incidencia legal frente a incidencia económica (o sujeto pasivo fente a repercutido).

\paragraph{Incidencia económica: La intención no garantiza el resultado (SIDGWICK, 1848)}
Conviene no idealizar esta formulación sin señalar su fragilidad, que en cierto modo se detectó casi de inmediato. 

SIDGWICK (1883), observó con precisión que la clasificación de MILL descansa por completo sobre una intención legislativa, no sobre ningún resultado económico verificable, y que no es siquiera un asunto sencillo enunciar con exactitud el principio general para determinar dónde recae realmente la carga de un impuesto, aun suponiendo conocidos todos los hechos relevantes. 

El problema de fondo, que SIDGWICK señala (pero no resuelve), es que la traslación que MILL da por supuesta depende de condiciones de mercado que la ley no puede garantizar por sí sola: un impuesto indirecto formalmente diseñado para trasladarse puede no trasladarse en absoluto si el mercado no lo permite, mientras que un impuesto formalmente directo puede, en determinadas condiciones, trasladarse igualmente. La clasificación de MILL describe con precisión la arquitectura jurídica que el legislador construye pero no puede garantizar que esa arquitectura produzca el resultado económico que se propone.

\textcolor{yellow}{\textbf{OJO -> Esto en realidad no es del todo correcto.} Que la Ley identifique claramente el sujeto pasivo y exige la trazabilidad del impuesto mediante facturas, no dice nada sobre la incidencia económica de este. Dicho de otra forma, no sabemos el contrafactual en ausencia del impuesto por lo que realmente, no tenemos ni idea sobre la incidencia económica, más allá de estimaciones mediante cálculos sobre las elasticidades de los bienes. \textbf{Convendría verificar si lo que aquí se expone se está refiriendo a esto o no.}  

De la intención a la obligación jurídica: cómo el IVA moderno resuelve la fragilidad de Mill

Lo que sí conviene retener es cómo el diseño técnico del IVA moderno responde, sin saberlo explícitamente, a la propia fragilidad que SIDGWICK detectó en MILL. 

MILL hablaba de una mera expectativa. El legislador espera e intenta que el productor traslade el impuesto, pero no se lo exige jurídicamente, dejando ese resultado a merced de las condiciones de mercado que SIDGWICK identificaba. El IVA moderno, en cambio, no se conforma con esperar: convierte esa traslación en una obligación jurídica explícita (repercusión), exigiendo por ley que el sujeto pasivo consigne el impuesto de forma separada y visible en la factura que emite a su cliente, de modo que la intención de MILL deja de ser una simple expectativa económica sujeta a las condiciones del mercado y se convierte en un mandato normativo verificable y sancionable. Esta transformación —de la intención económica de 1848 a la obligación jurídica del siglo XX— es exactamente el terreno que exploramos en el resto de este bloque: cómo, a lo largo de toda una cadena de producción, el ordenamiento jurídico moderno construye, transacción por transacción, la escisión entre quien paga a la Hacienda Pública y quien, en última instancia, soporta la carga.
}

\subsection{Intermediaciones: ¿qué son, para el impuesto, los diferentes eslabones de la cadena?}
Recuperemos el ejemplo numérico que desarrollamos, Un fabricante, mayorista y minorista, cada uno ingresando solo la diferencia entre el impuesto repercutido y el soportado. Veamos que sucede cuando intentamos responder a la pregunta. ¿Quién es, en sentido económico, el verdadero sujeto gravado en cada una de esas tres empresas? La respuesta, si la miras con detenimiento, es que ninguna de las dos primeras lo es en absoluto. 

El mayorista de nuestro ejemplo repercutió 30 y soportó 20: su posición neta, en términos de flujo de caja propio, es exactamente la misma que si el impuesto no existiera —recibe de su cliente el precio más el impuesto, paga a su proveedor el precio más el impuesto, e ingresa a la Hacienda Pública solo la diferencia, sin que ni un solo euro de esa diferencia salga, en términos reales, de su propio patrimonio. Lo mismo ocurre con el minorista. La única empresa cuya posición se ve alterada de forma definitiva e irreversible es la última de la cadena, porque su cliente no tiene, a su vez, ningún impuesto que deducir: el importe que ese consumidor paga de más, por el impuesto incluido en el precio, no vuelve nunca a él. Esto significa que, pese a que las tres empresas de la cadena son sujetos pasivos del impuesto, solo la relación entre el último eslabón y el consumidor final constituye una verdadera imposición económica. Todo lo que ocurre entre las empresas intermedias es, en rigor, un mecanismo de recaudación fraccionada y anticipada de un impuesto cuyo destinatario económico último es siempre el mismo, con independencia de cuántas manos intermedias haya atravesado el bien antes de llegar a él.

\subsubsection{Costes de intervenir: ¿está libre de incidencia el intermediario?}
Conviene, sin embargo, no simplificar en exceso esta conclusión. Afirmar que las empresas intermedias son meras recaudadoras podría sugerir, equivocadamente, que su papel en el sistema es gratuito o neutral en todos los sentidos, y no lo es. 

Aunque su posición fiscal neta sea nula, el cumplimiento de esa función recaudadora impone costes reales que sí recaen sobre la propia empresa y no se trasladan a nadie más. Los informes comparados sobre coste de cumplimiento tributario, como el estudio conjunto \textit{Paying Taxes} del Banco Mundial y PwC, que mide sistemáticamente el número de horas anuales que una empresa mediana tipo dedica a cumplir con sus obligaciones fiscales en distintos países del mundo, sitúan la gestión del IVA como uno de los componentes que más tiempo consume del conjunto de la carga administrativa tributaria de cualquier empresa, con diferencias muy marcadas entre países según la complejidad de sus regímenes de tipos diferenciados y sus exigencias documentales.\footnote{%
    Este coste es exactamente el que ya encontramos en la protesta de POUJADE y su Unión de Defensa de Comerciantes y Artesanos, contemporánea del propio nacimiento del IVA de LAURÉ, no se dirigía únicamente contra la acumulación piramidal del viejo impuesto en cascada, sino también, y de forma muy explícita, contra la percepción de que el pequeño comerciante francés estaba siendo reclutado, sin compensación alguna, como agente recaudador del Estado. Obligado a llevar una contabilidad más compleja, a asumir responsabilidad legal por errores de cálculo, y a adelantar de su propia tesorería el impuesto repercutido hasta el momento de la liquidación periódica, todo ello sin retribución alguna por ese servicio prestado a la Hacienda Pública. Es un resentimiento que sigue teniendo eco cada vez que un colectivo empresarial protesta contra nuevas obligaciones de facturación electrónica o de comunicación en tiempo real de sus operaciones.
}


\paragraph{Proveedor considerado: IVA en la economía digital}

\textcolor{red}{\textbf{OJO -> No queda claro el nexo}}

Cerramos con la manifestación más reciente de la extensión de la condición de recaudador a un tipo de actor económico particular: la plataforma digital de comercio electrónico. 

Desde 2021, el paquete de reforma del IVA para el comercio electrónico de la Unión introdujo la figura del proveedor considerado. Esto implica que cuando una plataforma, como Amazon o eBay, facilita determinadas ventas transfronterizas realizadas por vendedores terceros que operan a través de ella, la ley considera que es la propia plataforma quien realiza la venta al consumidor final a efectos de IVA, con independencia de que la propiedad real del bien nunca pase por sus manos. De esta forma, la operación se desdobla a efectos fiscales en: (i) una entrega del vendedor real a la plataforma, exenta; y, (ii) una entrega posterior de la plataforma al consumidor, sobre la que la plataforma debe repercutir, declarar e ingresar el IVA correspondiente.\footnote{%
    La motivación de esta reforma, dirigida a recuperar una estimación de fraude que la propia Comisión cifraba en torno a $5.000-7.000$ millones de euros anuales en el comercio electrónico transfronterizo. Cuando el sujeto pasivo tradicional resulta poco fiable como recaudador, el legislador no abandona la lógica de la recaudación intermedia, sino que la traslada a un actor distinto y más controlable, exactamente la misma decisión de diseño que ya vimos con otros protagonistas, desde la Roma de los portoria hasta el propio LAURÉ: gravar allí donde el Estado puede realmente ver y hacer cumplir la obligación, aunque quien efectivamente pague, en último término, siga siendo, como siempre, el consumidor final.
}

\subsubsection{Sujetos exentos: ¿es esta carga excesiva para algunos sujetos?}

\textcolor{red}{\textbf{OJO -> Reformular el apartado, no está del todo claro}}

Todo lo desarrollado presupone que cualquier empresa se incorpora al mecanismo de recaudación fraccionada del IVA. Sin embargo, la inmensa mayoría de los sistemas comparados introducen una excepción deliberada a esa lógica. 

Bajo lo que se conoce como umbral de registro, expresado normalmente como cifra de negocios anual, un negocio/empresario que se sitúe por debajo de dicho umbral queda completamente fuera del sistema: no repercute IVA en sus ventas, pero tampoco puede deducir el que soporta en sus compras. 

La razón de fondo es la misma que ya identificamos al hablar del coste de cumplimiento. Buena parte del coste administrativo de gestionar el IVA es, en gran medida, un coste fijo por empresa, casi independiente de su volumen de facturación. Para una gran empresa, ese coste fijo se diluye sobre una base de ingresos enorme y resulta insignificante; para un autónomo con una facturación mínima, ese mismo coste fijo puede representar una fracción sustancial de todo lo que gana, sin que la recaudación adicional que el Estado obtendría de gravarle compense, ni de lejos, el coste de gestión que le impondría.

\paragraph{Divergencias internacionales: ¿cuánto cuestan estas desviaciones?}
La cuantía exacta de este umbral varía de forma extraordinaria entre países, revelando que no existe ningún consenso técnico sobre dónde trazar la línea. 

Por ejemplo, el Reino Unido mantiene el umbral de registro más alto de toda la OCDE ($90.000$ libras anuales), junto con Suiza, muy por encima de la media de la Unión Europea (en torno a $44.000$ libras) y de la media de la OCDE en su conjunto (en torno a $34.700$ libras), dejando fuera del sistema a una cifra estimada de $3,2$ millones de pequeñas empresas británicas. España, en el extremo opuesto, no contempla ningún umbral general de exención: cualquier actividad empresarial o profesional queda sujeta al IVA desde el primer euro de facturación, con independencia de su volumen.\footnote{%
    Se trata de una de las pocas grandes economías europeas que ha optado deliberadamente por no establecer ningún umbral, pese a que la propia normativa comunitaria lo permite
} 

Un detalle técnico del propio diseño británico es que el umbral de baja se fija deliberadamente más bajo que el umbral de alta (88.000 libras frente a 90.000), precisamente para evitar que una empresa cuya facturación oscile año a año alrededor de la frontera tenga que entrar y salir del sistema de forma continua, con el coste administrativo que ello generaría tanto para la empresa como para la Hacienda Pública —una pequeña "zona de histéresis" que confirma, por sí sola, cuánto le preocupa al legislador la inestabilidad alrededor de esa frontera—. Esa preocupación no es infundada, y conecta directamente con el problema de agrupamiento o *bunching* que ya trabajamos en otros contextos de este temario. La razón es que la inmensa mayoría de los umbrales de IVA se diseñan como una muesca (*notch*) en sentido estricto y no como un tramo progresivo: superar el umbral en un solo euro no obliga a tributar únicamente sobre ese exceso, sino que obliga a repercutir IVA sobre la totalidad de la facturación desde el primer euro del ejercicio —un salto discontinuo y potencialmente muy costoso para el empresario que se sitúa justo por encima de la línea—, lo que genera un incentivo extraordinariamente concentrado a permanecer, deliberadamente, un euro por debajo.

LIU et al. (2021), con acceso a registros administrativos vinculados de IVA y de impuesto sobre sociedades de empresas británicas entre 2004 y 2014, documentan un fenómeno que a primera vista parece contradictorio. Por un lado, confirman una acumulación masiva de empresas justo por debajo del umbral de registro, exactamente el patrón de bunching ya esperado. Pero, por otro lado, documentan que casi la mitad de las empresas cuya facturación se sitúa por debajo del umbral se registran voluntariamente de todos modos, pese a no estar legalmente obligadas a hacerlo; dos comportamientos, en apariencia, mutuamente incompatibles dentro del mismo colectivo de pequeñas empresas.

Los autores resuelven esta aparente contradicción identificando qué características de la empresa predicen uno u otro comportamiento, con resultados exactamente opuestos entre sí. La inscripción voluntaria resulta más probable: (i) cuanto mayor es la proporción de insumos intermedios comprados a otras empresas; (ii) cuanto menor es la competencia en su mercado, y, (iii) cuanto menor es la proporción de sus ventas dirigida directamente al consumidor final. El bunching muestra exactamente el patrón inverso, pues es más frecuente: (i) cuanto menor es la proporción de insumos comprados a otras empresas; y, (ii) cuanto mayor es la proporción de ventas dirigida al consumidor final. 

La lógica económica que explica este resultado resulta muy intuitiva. Una empresa que vende principalmente a otras empresas registradas tiene un incentivo real a registrarse voluntariamente, porque sus clientes empresariales quieren recibir una factura con IVA deducible. Pensemoslo, una empresa no registrada resulta, para un cliente empresarial, un proveedor menos atractivo que un competidor registrado, precisamente porque el cliente no puede deducir nada de lo que le compra a quien está fuera del sistema. Por otro lado, una empresa que vende principalmente al consumidor final, que nunca deduce nada en absoluto, no tiene ese mismo incentivo, y para ella permanecer fuera del sistema, evitando repercutir IVA y, potencialmente, pudiendo ofrecer un precio final más bajo que un competidor registrado, resulta la estrategia racional. Los propios autores añaden una precisión, parte del bunching podría deberse no a una restricción real de la actividad económica de la empresa, sino a una simple infradeclaración de ventas para mantenerse artificialmente por debajo del umbral. Una distinción entre respuesta real y respuesta de mero cumplimiento que resulta, en la práctica, muy difícil de separar con los datos administrativos disponibles, pero que apunta a un problema de control adicional que cualquier diseño de umbral debe tener presente.


\subsection{Inversión del sujeto pasivo: ¿qué hacer ante el fraude del operador desaparecido?}

Los datos más recientes de la Comisión Europea sobre el llamado VAT Gap lo sitúan en $128.000$ millones de euros en 2023, el $9,5\%$ de toda la recaudación de IVA teóricamente exigible en el conjunto de la Unión. De esa cifra agregada, la parte específicamente atribuible al fraude de operador desaparecido intracomunitario se estima entre $12.500$ y $32.800$ millones de euros anuales, de forma sostenida durante todo el periodo 2010-2023. La Fiscalía Europea (EPPO) advirtió en 2025 de que estas cifras oficiales podrían infraestimar sustancialmente la magnitud real del fraude organizado: solo sus investigaciones activas a finales de 2025 apuntaban a pérdidas estimadas de $45.000$ millones de euros, con el fraude de ingresos representando apenas el $27\%$ de sus casos pero el $67\%$ del daño financiero total investigado.

\textcolor{\textbf{OJO -> Hay que hacer una breve introducción: \textit{parece perfecto pero no lo es...}}}

En su forma más simple, el esquema exige solo tres empresas y aprovecha exactamente el régimen de destino que ya presentamos: una empresa $A$, situada en un Estado miembro, vende mercancía a una empresa $B$, situada en otro Estado miembro, como entrega intracomunitaria exenta. La empresa $B$ adquiere así la mercancía sin soportar IVA alguno. A continuación, $B$ la revende dentro de su propio país a una empresa $C$, repercutiéndole el IVA doméstico correspondiente. $B$ cobra ese IVA de $C$ y desaparece sin ingresarlo jamás a la Hacienda Pública, dejando de existir como sujeto localizable, sin domicilio, sin libros contables, sin nada que perseguir. $C$, mientras tanto, es una empresa perfectamente inocente que ha comprado con una factura en apariencia intachable, y ejerce con normalidad su derecho a deducir el IVA que le repercutió $B$.

La exposición más deliberada y sofisticada de este esquema es lo que se conoce como fraude carrusel. Imaginemos que una empresa conducto $A$ realiza una entrega intracomunitaria exenta a un operador desaparecido $B$ en otro Estado miembro; $B$ adquiere la mercancía sin pagar IVA y realiza después una entrega doméstica a una tercera empresa $C$, el llamado bróker. El operador desaparecido cobra el IVA de su venta al bróker pero no lo ingresa, y desaparece. El bróker C solicita entonces la devolución del IVA que soportó en su compra a $B$. De este modo, la pérdida no es solo el IVA no ingresado por $B$, sino también el reembolso que el propio Estado paga a $C$. Pero hay más, la empresa $C$ puede, a su vez, realizar una nueva entrega intracomunitaria exenta a $A$, que puede volver a entregarla, exenta, a $B$, reiniciando exactamente el mismo patrón. Es esa repetición en círculo lo que explica el término carrusel.\footnote{%
    Para dificultar la investigación, la propia Comisión señala que el tramo entre $B$ y $C$ suele poblarse, además, de una o varias empresas intermedias llamadas colchón o pantalla (buffers): compañías que no participan del fraude, ingresan correctamente el IVA de sus propias ventas, pero cuya única función real es alargar y camuflar la cadena entre el operador desaparecido y el bróker, dificultando que la Administración reconstruya el rastro.
} La propia autoridad tributaria del Reino Unido resume la diferencia con una frase que conviene retener por su precisión: \textit{el fraude carrusel es como el fraude de adquisición simple, salvo que la mercancía no termina en manos de un consumidor final, en su lugar, da la vuelta, normalmente hasta regresar al país de origen}.

En definitiva, el fraude simple de operador desaparecido genera, como mucho, una pérdida equivalente al IVA de una sola transacción; el fraude carrusel, al repetir el ciclo varias veces con la misma mercancía (o sin que la mercancía llegue siquiera a moverse físicamente, verificándose la operación solo sobre el papel), puede multiplicar esa pérdida tantas veces como vueltas complete el carrusel antes de ser detectado, lo que explica por qué es, con diferencia, la variante más dañina y la que concentra el mayor esfuerzo de detección de las autoridades europeas.

\subsubsection{Erosión de la base: ¿qué mercancias se ven afectadas por este fenómeno?}
Inicialmente, este mecanismo se aplicó a un patrón de fraude muy concreto. Las mercancías elegidas por las redes de fraude comparten tres características: alto valor unitario, peso o volumen físico reducido y una demanda de mercado suficientemente líquida como para revenderlas sin dificultad. Los teléfonos móviles y los microchips fueron los productos estrella de este tipo de fraude en toda Europa durante los años 2000, precisamente por reunir esas tres condiciones a la perfección. Sin embargo, se ha ido ampliando progresivamente siendo de muy difícil solución.

\paragraph{Bienes intangibles: Derechos de Emisión de Gases de Efecto Invernadero (EU-ETS)}
Uno de los casos más dramáticos y mejor documentados es el de los derechos de emisión de gases de efecto invernadero del Régimen de Comercio de Derechos de Emisión de la Unión Europea (EU-ETS). Entre 2008 y 2009, organizaciones criminales descubrieron que estos derechos de emisión eran el vehículo de fraude carrusel más eficiente jamás disponible. ¿Por qué? 

Un operador compraba derechos de emisión en un Estado miembro sin soportar IVA (adquisición intracomunitaria exenta), los revendía de inmediato en otro Estado miembro repercutiendo el IVA correspondiente, y desaparecía sin ingresarlo, en un ciclo que podía repetirse varias veces en el mismo día, dada la instantaneidad de la transferencia electrónica de un activo sin existencia física. 

Europol, que abrió una investigación específica a finales de 2008 ante el inusitado incremento del volumen de negociación en varios mercados de carbono europeos, llegó a estimar que hasta el $90\%$ del volumen total de negociación en algunos de esos mercados correspondía a operaciones fraudulentas, con un coste total estimado para las arcas públicas europeas de en torno a $5.000$ millones de euros.\footnote{%
    La respuesta de los distintos Estados miembros, en ausencia inicial de una solución coordinada, fue dispar y ilustra cómo un mismo problema puede resolverse de formas muy distintas según la tradición administrativa de cada país. Francia y el Reino Unido optaron, en un primer momento, por la solución más drástica exonerar o aplicar un tipo cero a la negociación de derechos de emisión, renunciando lisa y llanamente a gravarla. Por otro lado, los Países Bajos y España adoptaron de inmediato la inversión del sujeto pasivo, trasladando la obligación de declarar el IVA al comprador desde el primer momento en que detectaron el patrón fraudulento. La Unión Europea armonizó finalmente esta última respuesta y autorizó a los Estados miembros, con carácter opcional y temporal, a aplicar la inversión del sujeto pasivo tanto a los derechos de emisión como a otros bienes ya identificados como especialmente vulnerables al fraude carrusel: microchips, teléfonos móviles, metales preciosos y, más sorprendentemente, incluso perfumes.
}

Constituye una paradoja irónica: el mismo mecanismo de crédito de factura que presentamos como la solución técnica elegante a la acumulación piramidal de la cascada es, exactamente por esa misma razón, la puerta que el crimen organizado ha aprendido a explotar.


\subsubsection{Inversión del sujeto: ¿?}
¿Cómo podemos solucionar este tipo de fraude? Si el fraude consiste en que un intermediario cobra el IVA de su cliente y desaparece sin ingresarlo, la solución más quirúrgica no es perseguir con más medios a ese intermediario después de que haya desaparecido, sino rediseñar la transacción para que nunca llegue a cobrar nada que pueda robar.

Ese es el mecanismo de la inversión del sujeto pasivo. En lugar de que el vendedor repercuta el IVA a su cliente y sea él quien deba ingresarlo, la propia ley traslada esa obligación al comprador, siempre que este sea un empresario registrado con derecho a deducción. El comprador se convierte, simultáneamente, en obligado a declarar el IVA devengado por la operación y en titular del derecho a deducirlo, en la misma autoliquidación y en el mismo momento —dos apuntes que se cancelan exactamente entre sí, sin que ni un solo euro cambie físicamente de manos entre comprador y vendedor por este concepto—. Si nunca hay flujo de caja de IVA entre las partes, no hay nada que el vendedor pueda cobrar y desaparecer con ello: tanto el fraude simple de operador desaparecido como su variante en carrusel quedan, para esa transacción concreta, estructuralmente imposibilitados.

En esencia, si nunca hay flujo de caja de IVA entre las partes, no hay nada que el vendedor pueda cobrar y desaparecer con ello: el fraude del operador desaparecido queda, para esa transacción concreta, estructuralmente imposibilitado.

\paragraph{Límites de esta solución: ¿por qué no generalizarla a toda la cadena?}
Si la inversión del sujeto pasivo elimina de raíz el fraude del operador desaparecido en cada transacción a la que se aplica, ¿por qué no generalizarla a todas las transacciones entre empresas, eliminando así el fraude carrusel por completo? La propuesta se debatió seriamente dentro de la Unión Europea entre 2016 y 2019, con la República Checa liderando la petición de una autorización especial para aplicarlo a escala nacional a todas las operaciones profesionales por encima de un determinado importe. 

La Comisión y la mayoría de los Estados miembros terminaron rechazando esa generalización, por dos razones. Primera, generalizar la inversión del sujeto pasivo no elimina el fraude, sino que lo desplaza por completo hacia la última fase de la cadena, concentrando en ese único punto todo el riesgo de fraude que antes se repartía, de forma diluida, a lo largo de toda la cadena de producción, y convirtiendo así cada operación minorista de cierto volumen en un objetivo de fraude potencialmente mucho más atractivo que cualquier eslabón intermedio aislado. 

Segunda, porque prescindir del mecanismo de crédito de factura en toda la cadena intermedia renunciaría precisamente a la propiedad de autocontrol documental que ya identificamos como una de las grandes virtudes del diseño de LAURÉ. Sin esa cadena de facturas cruzadas entre empresas independientes, cada una vigilando a la anterior por puro interés propio, el sistema perdería buena parte de su capacidad de verificación cruzada, sustituyendo un problema de fraude bien identificado y localizado por otro, potencialmente más difícil de detectar, disperso en el único punto de la cadena donde ya no hay ningún comprador posterior interesado en exigir una factura correcta.



\clearpage
\section{Dimensión distributiva: ¿?}
\puntosclave{
    
}

Queda, como última gran pregunta, la que probablemente más se discute en el debate público sobre la imposición indirecta y, paradójicamente, exige el aparato técnico más cuidadoso: ¿cómo se reparte, en la práctica, la carga de este impuesto entre los distintos niveles de renta de la población, y qué instrumentos tiene el legislador para modular ese reparto? 

El sentido común da por resuelta esta pregunta con una sola palabra: regresividad. Veamos como la respuesta es bastante más matizada recorriendo tres cuestiones: (i) cómo medir la regresividad y qué cambia cuando el horizonte de medición se alarga; (ii) si conviene diferenciar tipos según el bien; y, (iii) qué diferencia técnica separa la exención del tipo cero como instrumentos de corrección, con consecuencias que van mucho más allá de lo que su apariencia similar sugiere.

\subsection{Regresividad: ¿qué efectos redistributivos tiene la imposición indirecta?}
El argumento que cualquier manual introductorio de hacienda pública presenta como evidente es sencillo de enunciar: si se mide el impuesto pagado como proporción de la renta anual de cada hogar, el IVA resulta regresivo, porque la propensión a consumir decrece conforme aumenta la renta. 

La famosa utilidad marginal decreciente. Un hogar de renta baja gasta, típicamente, la totalidad de sus ingresos anuales, mientras que un hogar de renta alta ahorra una parte sustancial de la suya. Si ambos hogares soportan el mismo tipo de IVA sobre cada euro que gastan, el hogar de renta baja termina pagando, como proporción de su renta anual, un porcentaje sensiblemente mayor que el hogar de renta alta.

\subsubsection{Ciclo vital (POTERBA, 1989): ¿cambian las implicaciones distributivas incluyendo una dimensión temporal?}
Poterba (1989) plantea una objeción metodológica que, sin negar el argumento anterior, obliga a matizarlo de forma sustancial. Su punto de partida es una observación casi obvia una vez enunciada, pero con consecuencias de gran calado: la renta anual de un hogar es una medida muy ruidosa de su verdadera posición económica, porque está sujeta a perturbaciones transitorias que no reflejan su capacidad económica real y sostenida a lo largo del ciclo vital (cambios de empleo, jubilación, estudios, volatilidad, etc.). 

Por ejemplo, un jubilado que vive de sus ahorros acumulados durante toda una vida laboral puede tener una renta anual corriente muy baja, y sin embargo mantener un nivel de consumo notablemente superior a esa renta anual, financiado con el patrimonio ya acumulado. De la misma forma, un joven recién licenciado en su primer año de trabajo puede tener una renta anual también baja, pero un consumo elevado si anticipa que su renta futura será sustancialmente mayor y decide endeudarse o gastar sus ahorros del periodo de estudios en previsión de esa mejora. En ambos casos, la renta anual subestima de forma sistemática la verdadera capacidad económica a lo largo del ciclo vital.

POTERBA somete esta intuición a contraste empírico usando la Consumer Expenditure Survey estadounidense, centrándose en tres impuestos especiales (gasolina, alcohol y tabaco) y comparando dos formas de medir su regresividad: el gasto en esos bienes como proporción de la renta anual corriente, frente al gasto en esos mismos bienes como proporción del consumo total del hogar. Este segundo método siendo, según POTERBA, un proxy mucho más fiable de la renta permanente o vitalicia, precisamente porque las familias tienden a ajustar su nivel de consumo a su renta esperada de largo plazo, no a las fluctuaciones transitorias de su renta de un año concreto. Sus resultados son, en cierta medida, contundentes. El gasto en estos bienes, medido como proporción del consumo total, está mucho más igualmente distribuido entre hogares que el mismo gasto medido como proporción de la renta anual, lo que le lleva a concluir que, desde una perspectiva de horizonte más largo, estos impuestos especiales son sustancialmente menos regresivos de lo que sugieren los análisis convencionales basados en renta anual.

Aunque el trabajo original de POTERBA se centra en impuestos especiales sobre bienes concretos, su marco metodológico se trasladó directamente al IVA general apenas unos años después. CASPERSEN y METCALF (1994) aplican exactamente la misma lógica a un hipotético IVA general y confirman el mismo patrón cualitativo que POTERBA. Comparando la incidencia medida sobre renta anual frente a la incidencia medida sobre una aproximación a la renta vitalicia, la regresividad aparente de un IVA de tipo único, medida en términos de renta vitalicia o de consumo total, resulta sustancialmente menor que la que arroja el cálculo convencional basado en renta anual. De hecho, en algunas de sus especificaciones, la diferencia entre ambas medidas resulta lo bastante grande como para que la conclusión cualitativa cambie de forma apreciable, acercando el IVA a la proporcionalidad en lugar de a la regresividad marcada que sugiere el análisis anual.

\paragraph{Ahorro: Propensión creciente en renta}
Conviene aclarar, no obstante, que el hallazgo de POTERBA y de CASPERSEN y METCALF no significa que el IVA sea progresivo, ni siquiera que la preocupación distributiva convencional carezca de fundamento. Significa que la magnitud exacta de la regresividad depende de forma crítica de qué horizonte temporal y qué variable de referencia se elija para medirla. 

La objeción que cabe hacer, en sentido inverso, es que incluso midiendo sobre renta vitalicia o sobre consumo total, la propensión a ahorrar una parte de esa renta vitalicia puede seguir siendo sistemáticamente mayor entre los hogares de renta vitalicia más alta, de modo que una parte de regresividad genuina podría persistir incluso bajo esta óptica más exigente. 

La conclusión razonable no es que el debate esté resuelto en ningún sentido, sino que la elección del horizonte de medición no es neutral, y que cualquier afirmación tajante sobre la regresividad del IVA (o no) que no explicite sobre qué variable y qué horizonte se está calculando, debe leerse con cautela.

\subsection{Progresividad: ¿podemos resolver el carácter regresivo del impuesto?}

\textcolor{red}{\textbf{OJO -> Hacer una pequeña introducción}}


\subsubsection{Contradicción teórica: ¿deberíamos o no diferenciar por tipos?}

\textcolor{red}{\textbf{OJO -> Hacer una pequeña introducción}}

\paragraph{Teorema de ATKINSON y STIGLITZ (1976): No}
ATKINSON y STIGLITZ (1976) demuestran un resultado que sigue siendo referencia obligada respecto al IVA. 

Su teorema establece que, si el Estado ya dispone de un impuesto óptimo y no lineal sobre la renta, y si las preferencias de los individuos cumplen una condición técnica llamada separabilidad débil entre el consumo de bienes y el ocio:\footnote{%
    Es decir, si la proporción en que un individuo prefiere un bien de consumo frente a otro no depende de cuántas horas trabaja.
} diferenciar los tipos de IVA entre bienes no aporta ninguna ganancia de bienestar adicional, ya que toda la labor redistributiva óptima debe descansar sobre la tarifa del impuesto sobre la renta. 

De esta forma, la diferenciación de tipos de IVA resultaría una distorsión innecesaria que no mejora la equidad del sistema, porque el impuesto sobre la renta, correctamente diseñado, ya logra exactamente el mismo objetivo redistributivo.

Es difícil exagerar lo contraintuitivo que resulta este resultado frente a la práctica real de cualquier sistema fiscal comparado. Prácticamente todos los países del mundo que aplican IVA general tienen tipos reducidos para bienes de consumo esencial (alimentos, medicamentos, etc.) justificados con el argumento de proteger a los hogares de menor renta. La conclusión no deja de tener su propia ironía: o bien ningún gobierno del mundo ha seguido nunca la recomendación teórica más citada sobre esta materia, o bien la condición de separabilidad que sostiene el resultado no se cumple en la realidad, y esta segunda posibilidad es, precisamente, la que CORLETT y HAGUE habían anticipado veintitrés años antes.

\paragraph{Teorema de CORLETT y HAGUE (1953): Sí}
CORLETT y HAGUE (1953) habían llegado, mucho antes que ATKINSON y STIGLITZ, a una conclusión que apunta en la dirección opuesta. Dado que el ocio no puede gravarse directamente, la imposición sobre la renta del trabajo distorsiona ya, por sí sola, la elección entre trabajar y disfrutar de ocio, empujando a cualquier individuo racional hacia más ocio y menos trabajo del que elegiría en ausencia de ese impuesto. CORLETT y HAGUE demuestran que, dado esa distorsión de partida, resulta óptimo compensar parcialmente ese sesgo gravando con mayor intensidad los bienes que son complementarios del ocio, y con menor intensidad los bienes que son sustitutivos del ocio. Bajo esta lógica, diferenciar los tipos de IVA no es una distorsión innecesaria sino el instrumento que corrige, indirectamente, la propia distorsión que el impuesto sobre la renta introduce en la elección trabajo-ocio.

\paragraph{Reconciliación de KAPLOW (2008): Depende}
Durante décadas, ambos resultados convivieron en la literatura como una aparente contradicción sin resolver del todo.\footnote{
Incluso agravada por el hecho de que algunos trabajos posteriores, utilizando el marco ATKINSON-STIGLITZ y extendiendolo al análisis de CORLETT-HAGUE dentro de un Sistema Fiscal completo, con impuesto sobre la renta plenamente óptimo y no lineal, habían llegado a conclusiones que invertían el propio signo de la recomendación CORLETT y HAGUE.
} Sin embargo, KAPLOW (2008) aclaró definitivamente esta confusión, demostrando que la recomendación original de CORLETT y HAGUE sigue siendo correcta incluso bajo un marco de imposición óptima y no lineal sobre la renta, y que las derivaciones posteriores que habían llegado a la conclusión contraria contenían un error técnico identificable. 

Lo verdaderamente valioso de la aclaración de KAPLOW es la forma en que delimita con precisión el terreno de cada uno de los dos resultados, mostrando que no son, en realidad, contradictorios sino anidados. El resultado de uniformidad de ATKINSON y STIGLITZ es válido exactamente cuando se cumple la condición de separabilidad débil. Es decir, cuando ningún bien concreto guarda una relación especial con el ocio distinta de la de los demás. Por otro lado, la recomendación de diferenciación de CORLETT y HAGUE es la extensión correcta para el caso, mucho más realista, en que esa separabilidad falla, porque algunos bienes son, de hecho, complementarios o sustitutivos del ocio en un grado claramente distinto al de otros.

Esta delimitación tiene, además, una proyección directa sobre debates de política fiscal muy actuales. Varios países europeos han justificado, explícita o implícitamente, tipos reducidos de IVA sobre servicios como la restauración o los servicios domésticos de proximidad precisamente con un argumento de corte CORLETT-HAGUE: comer fuera de casa, o contratar ayuda doméstica, libera tiempo que de otro modo se destinaría a tareas del hogar, permitiendo a quien lo contrata dedicar más horas al trabajo remunerado; exactamente lo contrario de lo que recomendaría la uniformidad de ATKINSON y STIGLITZ. Evidentemente, ningún gobierno ha diseñado nunca su catálogo de tipos consultando ningún tipo de \textit{matriz de complementariedad} de cada bien con el ocio, pues resulta imposible de componer (a día de hoy), pero el patrón de tipos reducidos que la práctica comparada ha ido adoptando de forma más bien intuitiva y políticamente oportunista guarda un parecido razonable con lo que esa teoría recomendaría.

\subsubsection{Materialización práctica: ¿exención o tipo cero, cuál es mejor?}
Cerramos este apartado con una distinción técnica que determina si un beneficio fiscal aparentemente generoso cumple realmente su propósito o no por su propio diseño. 

Una venta exenta de IVA no repercute impuesto alguno al comprador, pero el vendedor exento no puede deducir el IVA que él mismo soportó en sus propias compras. ¿Por qué? Como es obvio, queda, a todos los efectos, fuera de la cadena (exactamente en la misma posición que el pequeño negocio no registrado). Por el contrario, una venta a tipo cero sigue siendo formalmente una operación sujeta a IVA (del $0\%$), lo que significa que el vendedor sí conserva el derecho pleno a deducir todo el IVA que soportó en sus compras, permaneciendo íntegramente dentro de la cadena de neutralidad que DIAMOND y MIRRLEES identificaron como la propiedad esencial de cualquier diseño óptimo de imposición sobre bienes intermedios.

\paragraph{Exención: Reversión al impuesto plurifásico}
Esta diferencia, que puede parecer un puro tecnicismo administrativo, tiene una consecuencia que conecta directamente con el problema central: la exención reintroduce, a pequeña escala y de forma invisible, el mismo problema de acumulación en cascada que el propio LAURÉ diseñó su sistema para eliminar. 

Imaginemos una empresa cuya actividad está exenta, pero que soporta $20\%$ de IVA en sus propias compras de bienes y servicios necesarios para prestar ese servicio. Como no puede deducir ese IVA, lo incorpora, como cualquier otro coste no recuperable, al precio que cobra a sus clientes. Si ese cliente es, a su vez, un consumidor final, el problema queda ahí: el consumidor paga un precio que incluye, de forma oculta y sin ningún desglose visible, ese IVA no deducido. Pero ojo, aunque esto no incumpla la optimalidad en términos de DIAMOND y MIRRLESS, sucede exactamente lo que la exención pretendía evitar, solo que camuflado dentro del precio en lugar de mostrado como impuesto explícito. 

Ahora bien, si ese cliente es otra empresa que compra el servicio exento como insumo de su propia actividad, el problema es otro. La segunda empresa no tiene ninguna factura con IVA desglosado que deducir, de modo que el IVA no deducido por la primera empresa queda enterrado dentro del precio que paga la segunda, sin ninguna posibilidad de recuperarlo, y se traslada así, silenciosamente, a la siguiente fase de la cadena, exactamente el mecanismo de acumulación piramidal que el propio LAURÉ quería evitar, incumpliendo la optimalidad en términos de DIAMOND y MIRRLESS.

El ejemplo que mejor ilustra por qué, pese a esta consecuencia indeseable, la exención sigue siendo la solución mayoritaria en la práctica comparada, es el de los servicios financieros. 

Prácticamente todos los sistemas de IVA del mundo eximen la práctica totalidad de los servicios financieros de intermediación, y la razón declarada no es de política social ni de mérito, sino puramente técnica. El precio que un banco cobra por intermediar entre depositantes y prestatarios no adopta la forma de una comisión explícita y fácilmente identificable transacción por transacción, sino la de un margen o diferencial entre el tipo de interés que paga a sus depositantes y el que cobra a sus prestatarios. Un margen que se conoce con precisión a nivel agregado de toda la cartera del banco, pero que resulta extraordinariamente difícil de descomponer, préstamo por préstamo y depósito por depósito, en la unidad de análisis que el crédito de factura exige.\footnote{%
    Hay excepciones. Por ejemplo, en España, mucha intermediación bancaria queda públicamente registrada, como escritura, por lo que a priori sí podría conocerse un extremo de la cadena (aunque no el otro). Por otro lado, algunos servicios financieros basados en comisiones explícitas, como la gestión de una cuenta o una transferencia con cargo fijo, sí podrían gravarse sin dificultad técnica alguna, y de hecho algunos países los gravan; es específicamente el componente de margen implícito el que ha resistido, durante más de medio siglo, cualquier intento de diseño técnico satisfactorio.
}

Consciente de que esta exención generaba exactamente la distorsión de cascada que acabamos de describir, el FMI propuso en 2010 una alternativa completamente distinta a intentar forzar el encaje del sector financiero dentro del propio IVA: el Impuesto sobre las Actividades Financieras. Esta figura separada gravaría directamente la suma de los beneficios y las remuneraciones del sector financiero.\footnote{%
    La propia Comisión había intentado, en 2007, modernizar el tratamiento del IVA financiero mediante dos propuestas legislativas específicas, que quedaron prácticamente paralizadas cuando estalló la crisis financiera ese mismo año.
}

\paragraph{Tipo cero: Preservar la neutralidad}
Frente a este uso forzado de la exención por pura imposibilidad técnica, el tipo cero se reserva, de forma consciente y deliberada, para los casos en que el legislador sí quiere preservar la plena neutralidad de la cadena, pero desea que el resultado final para el consumidor sea de carga nula.

Las exportaciones son el caso universal. Se trata de transacciones desgravadas a tipo cero, precisamente para que el exportador conserve el derecho a deducir todo su IVA soportado, evitando que el producto exportado cargue, silenciosamente, con impuesto español o europeo al competir en un mercado extranjero. 

Es cierto que algunos países han llevado esta misma lógica, más allá del comercio exterior, a determinados bienes de primera necesidad por razones explícitamente distributivas. Por ejemplo, el Reino Unido sobre la mayoría de los alimentos, los libros y la ropa infantil que se remonta a acuerdos previos a su propia incorporación a la Unión y que negoció mantener como cláusula \textit{stand-still} al adherirse en 1973. Esto contrasta con la práctica habitual en el resto de países de gravar estos bienes a tipos reducidos. Diferencia que no es meramente simbólica, ya que bajo tipo reducido el consumidor sigue pagando algo, mientras que bajo tipo cero no paga nada en absoluto por el propio impuesto, aunque el vendedor conserve, a diferencia de la exención, el derecho pleno a deducir cuanto haya soportado. Siendo esta última la combinación más generosa posible para el consumidor final que el diseño técnico del IVA permite, y la que menos países se permiten aplicar fuera del ámbito estricto del comercio exterior, precisamente por su coste recaudatorio.














































\clearpage
\section{Conclusión} 


\subsection{Recapitulación}


\subsection{Extensiones y relación con otras partes del Temario}



\subsection{Opinión}

\subsubsection{Idea final (Salida o cierra)}

\clearpage
\pagenumbering{gobble}
\MyBibliography

- HOBBES, T. (1651): *Leviathan, or the Matter, Forme and Power of a Commonwealth Ecclesiasticall and Civil*, capítulo 30, "Of the Office of the Sovereign Representative".
- RAWLS, J. (1971): *A Theory of Justice*, Harvard University Press, Cambridge (Massachusetts), sección sobre elección entre impuesto sobre la renta y sobre el gasto.
- SCHNEIDER, F. (2007): "The Size of the Shadow Economies of 145 Countries all over the World: First Results over the Period 1999 to 2003", *Journal of Population Economics*, vol. 20, núm. 3.
Sistema de portoria del Imperio romano (fuentes clásicas y epigráficas sobre aduanas provinciales).
Decreto de la Asamblea Nacional Constituyente francesa de 1790 (abolición de las aduanas internas, traites).
Historia institucional de la Ferme Générale y de la gabela de la sal en la Francia del Antiguo Régimen.
BENSAUDE-VINCENT, B. (1993): Lavoisier: Mémoires d'une révolution, Flammarion, París.
LAURÉ, M. (1953): Traité de politique fiscale, Presses Universitaires de France.
Primera y Segunda Directivas IVA del Consejo, de 11 de abril de 1967.
Tratado Constitutivo de la Comunidad Económica Europea (Tratado de Roma), 1957.
- MILL, J. S. (1848): *Principles of Political Economy*, Libro V, capítulo VI.
- SAUSGRUBER, R. y TYRAN, J.-R. (2005): "Testing the Mill Hypothesis of Fiscal Illusion", *Public Choice*, vol. 122, núms. 1-2.
- PUVIANI, A. (1903): *Teoria della illusione finanziaria*, Sandron, Palermo.
- SMITH, A. (1776): *An Inquiry into the Nature and Causes of the Wealth of Nations*, Libro V (cuarto canon de la imposición, "economía").
- CHETTY, R., LOONEY, A. y KROFT, K. (2009): "Salience and Taxation: Theory and Evidence", *American Economic Review*, vol. 99, núm. 4.
- Comisión Europea (2023): *VAT Gap Report*, Dirección General de Fiscalidad y Unión Aduanera.
- WIFO / Comisión Europea: *VAT Compliance Gap Due to Missing Trader Intra-Community (MTIC) Fraud*, estimaciones 2010-2023.
- Fiscalía Europea (EPPO) (2025): declaraciones ante el Parlamento Europeo sobre el alcance del fraude de IVA y aduanas, septiembre de 2025.
South Dakota v. Wayfair, Inc., 585 U.S. ___ (2018).
Quill Corp. v. North Dakota, 504 U.S. 298 (1992).
National Bellas Hess, Inc. v. Department of Revenue of Illinois, 386 U.S. 753 (1967).
DIAMOND, P. A. y MIRRLEES, J. A. (1971): "Optimal Taxation and Public Production I: Production Efficiency" y "II: Tax Rules", American Economic Review, vol. 61, núms. 1 y 3.
Historia económica del Umsatzsteuer alemán y su relación con la concentración empresarial vertical en el periodo de entreguerras.
- Agencia Nacional Tributaria de Japón: introducción del Sistema de Factura Cualificada (*inboisu seido*), 1 de octubre de 2023.
- Ley de reforma del Impuesto sobre el Consumo japonés de 2019 (introducción del tipo reducido del 8%).
- Comisión Europea (2024): paquete legislativo *VAT in the Digital Age* (ViDA).
- *Sistema di Interscambio* (SdI), sistema italiano de facturación electrónica obligatoria, vigente desde 2019.
- Comisión Europea (1985): *Libro Blanco para la Realización del Mercado Interior*, COM(85) 310.
- Directiva 91/680/CEE del Consejo, de 16 de diciembre de 1991 (régimen transitorio de IVA para las transacciones intracomunitarias, vigente desde el 1 de enero de 1993).
- Comisión Europea (2019): propuestas legislativas hacia un "sistema definitivo de IVA" basado en el refuerzo del principio de destino.
- Comisión Europea (2025): *Programa de Trabajo de la Comisión para 2025*, confirmación de la retirada de la propuesta de sistema definitivo de IVA, 11 de febrero de 2025.
- Acuerdo General sobre Aranceles Aduaneros y Comercio (GATT) y Acuerdo de la OMC sobre Subvenciones y Medidas Compensatorias (asimetría entre ajuste en frontera de impuestos directos e indirectos).
- Propuesta de *Destination-Based Cash Flow Tax* (DBCFT), Cámara de Representantes de los Estados Unidos, 2016-2017 ("A Better Way", Paul Ryan y Kevin Brady).
- PIGOU, A. C. (1920): *The Economics of Welfare*, Macmillan, Londres.
- COLCHERO, M. A. et al. (2016): "Beverage Purchases from Stores in Mexico under the Excise Tax on Sugar Sweetened Beverages: Observational Study", *BMJ*, vol. 352.
- COLCHERO, M. A. et al. (2017): "In Mexico, Evidence of Sustained Consumer Response Two Years after Implementing a Sugar-Sweetened Beverage Tax", *Health Affairs*, vol. 36, núm. 3.
- RAMSEY, F. P. (1927): "A Contribution to the Theory of Taxation", *The Economic Journal*, vol. 37, núm. 145 (ya citado; ahora en su formulación original sobre bienes de consumo).
- SUITS, D. B. y MUSGRAVE, R. A. (1953): "Ad Valorem and Unit Taxes Compared", *Quarterly Journal of Economics*, vol. 67, núm. 4.
- MILL, J. S. (1848): *Principles of Political Economy*, Libro V, capítulo III, "Of Direct Taxes", §1.
- SIDGWICK, H. (1883): *The Principles of Political Economy*, Libro III, capítulo VIII, "Public Finance", §8.
- SISMONDI, J. C. L. de (1815): *Political Economy*, sobre el origen fisiócrata de la terminología "impuesto directo/indirecto" en la escuela de Quesnay.
- Reglamento de Ejecución (UE) 2019/2026 del Consejo, de 21 de noviembre de 2019 (normas de aplicación del régimen de "proveedor considerado" para plataformas digitales).
- Directiva (UE) 2017/2455 del Consejo, de 5 de diciembre de 2017 (paquete de reforma del IVA para el comercio electrónico).
- Banco Mundial y PwC: *Paying Taxes*, informe comparado anual sobre coste de cumplimiento tributario empresarial (ediciones sucesivas).
- Comisión Europea: estimaciones de fraude de IVA en el comercio electrónico transfronterizo, paquete de reforma de 2021.
- LIU, L., LOCKWOOD, B., ALMUNIA, M. y TAM, E. H. F. (2021): "VAT Notches, Voluntary Registration, and Bunching: Theory and UK Evidence", *The Review of Economics and Statistics*, vol. 103, núm. 1.
- HM Revenue & Customs (Reino Unido): *VAT registration threshold*, actualización de abril de 2024 (Value Added Tax Act 1994, Anexos 1 y 9ZA).
- Directiva (UE) 2020/285 del Consejo, de 18 de febrero de 2020 (régimen especial de las pequeñas empresas a efectos del IVA, umbrales máximos permitidos a los Estados miembros).
- Directiva 2010/23/UE del Consejo, de 16 de marzo de 2010, por la que se modifica la Directiva 2006/112/CE relativa al sistema común del IVA, en lo que respecta a la aplicación optativa y temporal del mecanismo de inversión del sujeto pasivo.
- Europol (2009): investigación sobre fraude organizado en los mercados europeos de derechos de emisión, diciembre de 2009.
- *Value Added Tax (Emissions Allowances) Order 2009* (SI 2009/2093), Reino Unido.
- Comisión Europea: propuesta de mecanismo generalizado de inversión del sujeto pasivo, debate 2016-2019 (solicitud de excepción de la República Checa).
POTERBA, J. M. (1989): "Lifetime Incidence and the Distributional Burden of Excise Taxes", American Economic Review, vol. 79, núm. 2.
CASPERSEN, E. y METCALF, G. E. (1994): "Is a Value Added Tax Regressive? Annual Versus Lifetime Incidence Measures", National Tax Journal, vol. 47, núm. 4.
LYON, A. B. y SCHWAB, R. M. (1991): "Consumption Taxes in a Life-Cycle Framework: Are Sin Taxes Regressive?", NBER Working Paper 3932.
- ATKINSON, A. B. y STIGLITZ, J. E. (1976): "The Design of Tax Structure: Direct versus Indirect Taxation", *Journal of Public Economics*, vol. 6, núms. 1-2.
- CORLETT, W. J. y HAGUE, D. C. (1953): "Complementarity and the Excess Burden of Taxation", *The Review of Economic Studies*, vol. 21, núm. 1.
- KAPLOW, L. (2008): "Taxing Leisure Complements", NBER Working Paper 14397.
- Directiva 2006/112/CE del Consejo, de 28 de noviembre de 2006, relativa al sistema común del Impuesto sobre el Valor Añadido, arts. 135.1 y 137 (exención de servicios financieros y opción de gravamen).
- Fondo Monetario Internacional (2010): *A Fair and Substantial Contribution by the Financial Sector*, informe final para el G20 (propuesta del Financial Activities Tax).
- PODDAR, S. (2010): "Consumption Taxes: The Role of the Value-Added Tax", capítulo sobre fiscalidad del sector financiero.
- KERRIGAN, A. (2010): "The Elusiveness of Neutrality – Why Is It So Difficult To Apply VAT to Financial Services?", *International VAT Monitor*, vol. 21, núm. 2.
- Comisión Europea (2007): propuestas COM(2007)747 y COM(2007)746 sobre modernización del régimen de IVA de servicios financieros y de seguros.
- Value Added Tax Act 1994 (Reino Unido), Anexo 8 (régimen de tipo cero sobre alimentos, libros y ropa infantil, cláusula de *standstill* de 1973).

\MyBibItem{DAWID y GATTI}{Chapter 2 - Agent-Based Macroeconomics}{2018}{https://doi.org/10.1016/bs.hescom.2018.02.006}


% ANEXOS
\clearpage










\end{document}